\subsection{\texorpdfstring{The complexes \(\mathsf{K}(u),\mathsf{K}.(u,\mathsf{C}),\mathsf{K}^{\prime}(u,\mathsf{C})\)}{The complexes \textbackslash mathsf\{K\}(u),\textbackslash mathsf\{K\}.(u,\textbackslash mathsf\{C\}),\textbackslash mathsf\{K\}\^{}\{\textbackslash prime\}(u,\textbackslash mathsf\{C\})}}

Let A be a ring, L an A-module, \(u: \mathbf{L} \to \mathbf{A}\) a linear form, and \(\Lambda(\mathbf{L})\) the exterior algebra of the A-module L. For \(x \in \Lambda(\mathbf{L})\)  let \(d_u(x)\) denote the interior product \(x \sqsubseteq u\) (III, p.~161, example). According to loc. cit., p.~162, formula (60), we have

\[
d _ {u} \left(e _ {1} \wedge \dots \wedge e _ {n}\right) = \sum_ {i = 1} ^ {n} (- 1) ^ {i + 1} u \left(e _ {i}\right) e _ {1} \wedge \dots \wedge e _ {i - 1} \wedge e _ {i + 1} \wedge \dots \wedge e _ {n}\tag{1}
\]

for \(e_1, \ldots, e_n\) in L. According to III, p.~164 and 165, the map \(d_u: \symbf{\Lambda}(\mathbf{L}) \to \symbf{\Lambda}(\mathbf{L})\) is an antiderivation of degree \((-1)\) and of square zero. It is the unique antiderivation of the A-algebra \(\symbf{\Lambda}(\mathbf{L})\) which extends \(u: \symbf{\Lambda}^1(\mathbf{L}) \to \symbf{\Lambda}^0(\mathbf{L})\) .

DEFINITION 1. --- The complex \((\symbf{\Lambda}(\mathbf{L}), d_u)\) is denoted \(\mathbf{K}^{\mathbf{A}}(u)\) or \(\mathbf{K}(u)\) .

Be careful that \(\mathbf{K}_{n}(u)=\symbf{\Lambda}^{n}(\mathbf{L})=\mathbf{K}^{-n}(u)\) It is clear that \(\mathbf{K}(u)\) is zero on the right and that \(\mathrm{H}_{0}(\mathbf{K}(u))=\mathrm{Coker}(u)=\mathrm{A}/\mathfrak{q}\) where \(\mathfrak{q}\) is the ideal \(u(\mathbf{L})\) of A.

For every complex C of A-modules, we set

\[
\mathbf {K} _ {\cdot} ^ {\mathrm{A}} (u, \mathrm{C}) = \mathrm{C} \otimes_ {\mathrm{A}} \mathbf {K} ^ {\mathrm{A}} (u), \quad \mathbf {K} _ {\cdot} ^ {\cdot} (u, \mathrm{C}) = \operatorname{Homgr} _ {\mathrm{A}} (\mathbf {K} ^ {\mathrm{A}} (u), \mathrm{C}).
\]

\[
\mathrm{H} ^ {\mathsf {A}} (u, \mathrm{C}) = \mathrm{H} (\mathrm{C} \otimes_ {\mathsf {A}} \mathbf {K} ^ {\mathsf {A}} (u)) \quad \mathrm{H} _ {\mathsf {A}} ^ {\bullet} (u, \mathrm{C}) = \mathrm{H} (\mathrm{Homgr} _ {\mathsf {A}} (\mathbf {K} ^ {\mathsf {A}} (u), \mathrm{C})) ,
\]

\[
\mathrm{H} _ {r} ^ {\mathrm{A}} (u, \mathrm{C}) = \mathrm{H} _ {r} (\mathrm{C} \otimes_ {\mathrm{A}} \mathbf {K} ^ {\mathrm{A}} (u)), \quad \mathrm{H} _ {\mathrm{A}} ^ {r} (u, \mathrm{C}) = \mathrm{H} ^ {r} (\operatorname{Homgr} _ {\mathrm{A}} (\mathbf {K} ^ {\mathrm{A}} (u), \mathrm{C})) .
\]

We therefore have canonical homomorphisms of A-modules (X, p.~62 and p.~82)

\[
\gamma_ {0}: \mathrm{H} _ {0} (\mathbf {C}) \otimes_ {\mathbf {A}} \mathrm{A} / \mathfrak {q} \rightarrow \mathrm{H} _ {0} ^ {\mathbf {A}} (u, \mathbf {C}),
\]

\[
\lambda^ {0}: \mathrm{H} _ {\mathrm{A}} ^ {0} (u, \mathrm{C}) \rightarrow \operatorname{Hom} _ {\mathrm{A}} (\mathrm{A} / \mathfrak {q}, \mathrm{H} ^ {0} (\mathrm{C})).
\]

Lemma 1. --- If the complex C is zero on the right (resp. on the left), then \(\mathbf{K}^{\mathbf{A}}(u,\mathbf{C})\) (resp. \(\mathbf{K}_{\mathbf{A}}^{*}(u,c)\) ) is zero on the right (resp. on the left), and \(\gamma_{0}\) (resp. \(\lambda^{0}\) ) is bijective.

This follows from X, p.~62, prop. 1 and p.~82, prop. 1.

PROPOSITION 1. --- Let \(x \in L\) ; denote by \(\mathbf{R}_{x}: y \mapsto x \wedge y\) the left multiplication by \(x\) in the algebra \(\boldsymbol{\Lambda}(\mathbf{L})\) . Then \(d_{u} \circ \mathbf{R}_{x} + \mathbf{R}_{x} \circ d_{u} = u(x)\) . \(1_{\boldsymbol{\Lambda}(\mathbf{L})} = u(x)_{\boldsymbol{\Lambda}(\mathbf{L})}\) .

Indeed, \((d_u \circ R_x + R_x \circ d_u)(y) = d_u(x \wedge y) + x \wedge d_u(y)\) ; since \(d_u\) is a derivation, \(d_u(x \wedge y) + x \wedge d_u(y) = d_u(x) \wedge y = u(x).y\) .

COROLLARY 1. --- If u is surjective, \(\mathbf{K}(u)\) is homotopic to zero (X, p.~34) as is \(\mathbf{K}_{\cdot}^{\mathrm{A}}(u,\mathrm{C})\) and \(\mathbf{K}_{\cdot}^{\cdot}(\mathbf{u},\mathbf{C})\) for every complex C.

Indeed, there exists \(x \in \mathbf{L}\) such that \(u(x) = 1\) . Then \(\mathbf{K}(u)\) is homotopic to zero by prop. 1, hence also \(\mathbf{K}^{\mathrm{A}}(u, \mathbf{C})\) (X, p.~64, prop. 3) and \(\mathbf{K}_{\mathrm{A}}^{*}(u, \mathbf{C})\) (X, p.~83, prop. 3).

COROLLARY 2. --- Let C be a complex, Ann (C) its annihilator. Then q + Ann (C) annihilates H \(^{A}\) (u, C) and H \(^{*}_{A}\) (u, C).

For every \(\lambda \in q\) , the homothety \(\lambda_{\mathbf{K}(u)}\) is homotopic to zero by the proposition, hence also \(1_{C} \otimes \lambda_{\mathbf{K}(u)}\) and Homgr \((\lambda_{\mathbf{K}(u)}, 1_{C})\) by X, p.~64, prop. 3 and X, p.~83, prop. 3; it follows that \(\lambda\) annihilates H.(u, C) and H'(u, C). If \(\lambda \in \text{Ann}(C)\) , then \(1_{\mathbf{K}(u)} \otimes \lambda_{C}\) and Homgr \((1_{\mathbf{K}(u)}, \lambda_{C})\) are zero.

Suppose L is projective (resp. \(\mathbf{K}(u)\) acyclic in degrees \(>0\) ). Then the complex \(\symbf{\Lambda}(\mathbf{L})\) is projective by III, p.~87, cor. 2 (resp. is a resolution of A/q); by X, p.~102 (resp. p.~100), we therefore have, for every A-module M, homomorphisms

\begin{enumerate}
\def\labelenumi{(\arabic{enumi})}
\setcounter{enumi}{1}
\tightlist
\item
\end{enumerate}

\[
\mathrm{H} _ {r} ^ {\mathbf {A}} (u, \mathrm{M}) \rightarrow \operatorname{Tor} _ {r} ^ {\mathbf {A}} (\mathrm{A} / \mathrm{q}, \mathrm{M}), \quad \operatorname{Ext} _ {\mathbf {A}} ^ {r} (\mathrm{A} / \mathrm{q}, \mathrm{M}) \rightarrow \mathrm{H} _ {\mathbf {A}} ^ {r} (u, \mathrm{M})\tag{3}
\]

\[
\operatorname{Tor} _ {r} ^ {\mathrm{A}} (\mathrm{A} / \mathrm{q}, \mathrm{M}) \rightarrow \mathrm{H} _ {r} ^ {\mathrm{A}} (u, \mathrm{M}), \quad \mathrm{H} ^ {r} (u, \mathrm{M}) \rightarrow \operatorname{Ext} _ {\mathrm{A}} ^ {r} (\mathrm{A} / \mathrm{q}, \mathrm{M}).
\]

If L is projective and \(\mathbf{K}(u)\) acyclic in degrees \(>0\) , the homomorphisms (2) and (3) above are bijective and inverses of each other (X, p.~102, prop. 1).

PROPOSITION 2. --- Let \((\mathbf{L}_{i})_{i\in I}\) be a family of A-modules, where the set I is finite and totally ordered. Let u be a linear form on \(\bigoplus_{i\in I}\mathbf{L}_{i}\) , \(u_{i}\) its restriction to \(\mathbf{L}_{i}\) .

The canonical isomorphism of A-algebras (III, p.~84)

\[
g: \bigotimes_{\substack{i\in I}}\symbf{\Lambda}(\mathrm{L}_{i})\to \symbf{\Lambda}(\bigoplus_{i\in I}\mathrm{L}_{i})
\]

is an isomorphism of the complex \(\otimes \mathbf{K}(u_i)\) (X, p.~63) on the complex \(\mathbf{K}(u)\) .

Indeed, according to X, p.~64, remark 4, the differential D of the complex \(\otimes_{i\in I}\mathbf{K}(u_i)\) is an antiderivation; the antiderivations \(d_u\) and \(g\circ D\circ g^{-1}\) of \(\boldsymbol{\Lambda}(\oplus L_i)\) coincide on \(\oplus L_i\) with the map \(x\mapsto u(x)\) . 1 of \(\oplus L_i\) to \(\boldsymbol{\Lambda}(\oplus L_i)\) , hence they are equal (III, p.~128, cor.).

Let C and C' be two complexes of A-modules. We have (X, p.~63 and p.~99) canonical isomorphisms of complexes

\[
\mathbf {C} \otimes_ {\mathbf {A}} \left(\mathbf {C} ^ {\prime} \otimes_ {\mathbf {A}} \mathbf {K} (u)\right)\rightarrow \left(\mathbf {C} \otimes_ {\mathbf {A}} \mathbf {C} ^ {\prime}\right) \otimes_ {\mathbf {A}} \mathbf {K} (u)
\]

\[
\operatorname{Homgr} _ {\mathbf {A}} \left(\mathrm{C} ^ {\prime}, \operatorname{Homgr} _ {\mathbf {A}} (\mathbf {K} (u), \mathrm{C})\right)\rightarrow \operatorname{Homgr} _ {\mathbf {A}} \left(\mathrm{C} ^ {\prime} \otimes_ {\mathbf {A}} \mathbf {K} (u), \mathrm{C}\right),
\]

that is, isomorphisms

\begin{enumerate}
\def\labelenumi{(\arabic{enumi})}
\setcounter{enumi}{3}
\tightlist
\item
\end{enumerate}

\[
\mathbf {C} \otimes_ {\mathbf {A}} \mathbf {K} ^ {\mathrm{A}} (u, \mathbf {C} ^ {\prime}) \rightarrow \mathbf {K} ^ {\mathrm{A}} (u, \mathbf {C} \otimes_ {\mathbf {A}} \mathbf {C} ^ {\prime})\tag{5}
\]

\[
\operatorname{Homgr} _ {\mathbf {A}} \left(\mathrm{C} ^ {\prime}, \mathbb {K} _ {\mathbf {A}} ^ {*} (u, \mathrm{C})\right)\rightarrow \operatorname{Homgr} _ {\mathbf {A}} \left(\mathbb {K} _ {\mathbf {A}} ^ {*} (u, \mathrm{C} ^ {\prime}), \mathrm{C}\right).
\]

In (4) and (5), take \(\mathbf{C}^{\prime} = \mathbf{K}(u^{\prime})\) , where \(u^{\prime}: L^{\prime} \to A\) is a linear form on an A-module \(L^{\prime}\) and note that \(\mathbf{K}^{\mathrm{A}}(u, \mathbf{K}(u^{\prime}))\) which is equal by definition to \(\mathbf{K}(u^{\prime}) \otimes_{\mathrm{A}} \mathbf{K}(u)\) identifies, according to prop. 2, with \(\mathbf{K}(u^{\prime} \oplus u)\) where \(u^{\prime} \oplus u: L^{\prime} \oplus L \to A\) is the linear form \((x^{\prime}, x) \mapsto u^{\prime}(x^{\prime}) + u(x)\) We then obtain isomorphisms of complexes

\begin{enumerate}
\def\labelenumi{(\arabic{enumi})}
\setcounter{enumi}{5}
\tightlist
\item
\end{enumerate}

\[
\mathbf {K} _ {\bullet} ^ {\mathrm{A}} (u ^ {\prime} \oplus u, \mathrm{C}) \rightarrow \mathbf {K} _ {\bullet} ^ {\mathrm{A}} (u, \mathbf {K} _ {\bullet} ^ {\mathrm{A}} (u ^ {\prime}, \mathrm{C}))\tag{7}
\]

\[
\mathbf {K} _ {\mathrm{A}} ^ {\bullet} (u ^ {\prime}, \mathbf {K} _ {\mathrm{A}} ^ {\bullet} (u, \mathrm{C})) \rightarrow \mathbf {K} _ {\mathrm{A}} ^ {\bullet} (u ^ {\prime} \oplus u, \mathrm{C}).
\]

By passing to homology, we deduce isomorphisms of A-modules

\[
\begin{array}{c c} \mathrm{H} _ {r} ^ {\mathbf {A}} (u ^ {\prime} \oplus u, \mathrm{C}) \to \mathrm{H} _ {r} ^ {\mathbf {A}} (u, \mathbf {K} _ {\bullet} ^ {\mathbf {A}} (u ^ {\prime}, \mathrm{C}))  , & r \in \mathbf {Z}  , \\ \mathrm{H} _ {\mathbf {A}} ^ {r} (u ^ {\prime}, \mathbf {K} _ {\mathbf {A}} ^ {\bullet} (u, \mathrm{C})) \to \mathrm{H} _ {\mathbf {A}} ^ {r} (u ^ {\prime} \oplus u, \mathrm{C})  , & r \in \mathbf {Z}  . \end{array}
\]

Finally, note that the homomorphism deduced from the product in the algebra \(\Lambda(L)\)

\[
m: \mathbf {K} ^ {\mathrm{A}} (u) \otimes_ {\mathrm{A}} \mathbf {K} ^ {\mathrm{A}} (u) \rightarrow \mathbf {K} ^ {\mathrm{A}} (u)
\]

is a morphism of complexes (since \(d_{u}\) is an antiderivation). Assuming L is free of rank \(n\) and composing with the morphism of complexes \(\mathbf{K}^{\mathrm{A}}(u)\to\symbf{\Lambda}^{n}\mathrm{L}(-n)\) which is the identity in degree \(n\) from this one deduces a morphism of complexes

\[
\chi : \mathbf {K} ^ {\mathrm{A}} (u) \otimes_ {\mathrm{A}} \mathbf {K} ^ {\mathrm{A}} (u) \rightarrow \boldsymbol {\Lambda} ^ {n} \mathrm{L} (- n);
\]

to this morphism there corresponds canonically, according to X, p.~99, prop. 12, a morphism of complexes

\[
\varphi : \mathbf {K} ^ {\mathrm{A}} (u) \rightarrow \operatorname{Homgr} _ {\mathrm{A}} \left(\mathbf {K} ^ {\mathrm{A}} (u), \boldsymbol {\Lambda} ^ {n} \mathrm{L} (- n)\right)
\]

which is bijective (III, p.~87, formula (20)). For every complex C, one deduces from this a composed isomorphism

\[
\begin{array}{r l} \mathbf {K} _ {\bullet} ^ {\mathrm{A}} (u, \mathrm{C}) & = \mathrm{C} \otimes_ {\mathrm{A}} \mathbf {K} ^ {\mathrm{A}} (u) \xrightarrow {1 \otimes_ {\Phi}} \mathrm{C} \otimes \operatorname{Homgr} _ {\mathrm{A}} \left(\mathbf {K} ^ {\mathrm{A}} (u), \boldsymbol {\Lambda} ^ {n} \mathrm{L} (- n)\right) \to \\ & \to \operatorname{Homgr} _ {\mathrm{A}} \left(\mathbf {K} ^ {\mathrm{A}} (u), \mathrm{C} \otimes_ {\mathrm{A}} \boldsymbol {\Lambda} ^ {n} \mathrm{L} (- n)\right) = \mathbf {K} _ {\mathrm{A}} ^ {\bullet} (u, \mathrm{C} \otimes_ {\mathrm{A}} \boldsymbol {\Lambda} ^ {n} \mathrm{L} (- n)). \end{array}
\]

By passing to homology, one therefore has canonical isomorphisms

\[
\mathrm{H} _ {r} ^ {\mathbf {A}} (u, \mathrm{C}) \rightarrow \mathrm{H} _ {\mathbf {A}} ^ {n - r} (u, \mathrm{C} \otimes_ {\mathbf {A}} \boldsymbol {\Lambda} ^ {n} \mathrm{L}), \quad r \in \mathbf {Z}.\tag{8}
\]

Remarks. --- 1) * The preceding remains valid when L is projective of rank n.~*

\begin{enumerate}
\def\labelenumi{\arabic{enumi})}
\setcounter{enumi}{1}
\tightlist
\item
  Since L is free of rank n, \(\Lambda^{n}\) L is isomorphic to A, we have non-canonical isomorphisms \(\mathrm{H}_{r}^{\mathrm{A}}(u,\mathrm{C})\to\mathrm{H}_{\mathrm{A}}^{n-r}(u,\mathrm{C})\) .
\end{enumerate}

\subsection{Functoriality}

Let \(f: \mathbf{C} \to \mathbf{C}'\) be a morphism of complexes. We denote by

\[
\begin{array}{l} \mathbf {K} _ {\bullet} ^ {\mathrm{A}} (u, f): \mathbf {K} _ {\bullet} ^ {\mathrm{A}} (u, \mathrm{C}) \to \mathbf {K} _ {\bullet} ^ {\mathrm{A}} (u, \mathrm{C} ^ {\prime}), \\ \mathbf {K} _ {\mathrm{A}} ^ {\bullet} (u, f): \mathbf {K} _ {\mathrm{A}} ^ {\bullet} (u, \mathrm{C}) \to \mathbf {K} _ {\mathrm{A}} ^ {\bullet} (u, \mathrm{C} ^ {\prime}), \end{array}
\]

the morphisms of complexes \(f \otimes 1_{\mathbf{K}(u)}\) and \(\operatorname{Homgr}_{\mathbf{A}}(1_{\mathbf{K}(u)}, f)\) .

We denote by \(\mathrm{H}_{\bullet}^{\mathrm{A}}(u,f):\mathrm{H}_{\bullet}^{\mathrm{A}}(u,\mathbf{C})\to\mathrm{H}_{\bullet}^{\mathrm{A}}(u,\mathbf{C}^{\prime}),\quad\mathrm{H}_{\mathrm{A}}^{\bullet}(u,f):\mathrm{H}_{\mathrm{A}}^{\bullet}(u,\mathbf{C})\to\mathrm{H}_{\mathrm{A}}^{\bullet}(u,\mathbf{C}^{\prime}),\quad\mathrm{H}_{r}^{\mathrm{A}}(u,f):\mathrm{H}_{r}^{\mathrm{A}}(u,\mathbf{C})\to\mathrm{H}_{r}^{\mathrm{A}}(u,\mathbf{C}^{\prime}),\mathrm{H}_{\mathrm{A}}^{r}(u,f):\mathrm{H}_{\mathrm{A}}^{r}(u,\mathbf{C})\to\mathrm{H}_{\mathrm{A}}^{r}(u,\mathbf{C}^{\prime})\) the induced morphisms in homology. The map \(f\mapsto\mathbf{K}_{\bullet}^{\mathrm{A}}(u,f)\) is linear; if \(g:\mathbf{C}^{\prime}\to\mathbf{C}^{\prime\prime}\) is another morphism of complexes, we have \(\mathbf{K}_{\bullet}^{\mathrm{A}}(u,g\circ f)=\mathbf{K}_{\bullet}^{\mathrm{A}}(u,g)\circ\mathbf{K}_{\bullet}^{\mathrm{A}}(u,f)\) ; similarly for \(\mathbf{K}_{\mathrm{A}}^{\bullet},\mathrm{H}_{\bullet}^{\mathrm{A}},\mathrm{H}_{\mathrm{A}}^{\bullet},\mathrm{H}_{r}^{\mathrm{A}},\mathrm{H}_{\mathrm{A}}^{r}\) .

Let \(0 \to \mathbf{C}' \xrightarrow{f} \mathbf{C} \xrightarrow{g} \dot{\mathbf{C}}'' \to 0\) be an exact sequence of complexes.

\begin{enumerate}
\def\labelenumi{\alph{enumi})}
\tightlist
\item
  Suppose L is flat; then \(\Lambda(L)\) is flat (X, p.~15, cor.). The sequence
\end{enumerate}

\[
0 \rightarrow \mathbf {K} _ {\bullet} ^ {\mathrm{A}} (u, C ^ {\prime}) \xrightarrow {\mathbf {K} _ {\bullet} ^ {\mathrm{A}} (u , f)} \mathbf {K} _ {\bullet} ^ {\mathrm{A}} (u, C) \xrightarrow {\mathbf {K} _ {\bullet} ^ {\mathrm{A}} (u , g)} \mathbf {K} _ {\bullet} ^ {\mathrm{A}} (u, C ^ {\prime \prime}) \rightarrow 0
\]

is then exact, and gives rise (X, p.~30) to an exact homology sequence

\[
\dots \to \mathrm{H} _ {n} ^ {\mathbf {A}} (u, \mathrm{C} ^ {\prime}) \xrightarrow {\mathrm{H} _ {n} (u , f)} \mathrm{H} _ {n} ^ {\mathbf {A}} (u, \mathrm{C}) \xrightarrow {\mathrm{H} _ {n} (u , g)} \mathrm{H} _ {n} ^ {\mathbf {A}} (u, \mathrm{C} ^ {\prime \prime}) \xrightarrow {\hat {\sigma} _ {n}} \mathrm{H} _ {n - 1} ^ {\mathbf {A}} (u, \mathrm{C} ^ {\prime}) \to \dots .
\]

\begin{enumerate}
\def\labelenumi{\alph{enumi})}
\setcounter{enumi}{1}
\tightlist
\item
  Suppose L is projective; then \(\Lambda(L)\) is projective. The sequence
\end{enumerate}

\[
0 \rightarrow \mathbf {K} _ {\mathrm{A}} ^ {*} (u, \mathrm{C} ^ {\prime}) \xrightarrow {\mathbf {K} _ {\mathrm{A}} ^ {*} (u , f)} \mathbf {K} _ {\mathrm{A}} ^ {*} (u, \mathrm{C}) \xrightarrow {\mathbf {K} _ {\mathrm{A}} ^ {*} (u , g)} \mathbf {K} _ {\mathrm{A}} ^ {*} (u, \mathrm{C} ^ {\prime \prime}) \rightarrow 0
\]

is then exact, and gives rise to an exact homology sequence

\[
\dots \rightarrow \mathrm{H} _ {\mathrm{A}} ^ {n} (u, \mathrm{C} ^ {\prime}) \xrightarrow {\mathrm{H} ^ {n} (u , f)} \mathrm{H} _ {\mathrm{A}} ^ {n} (u, \mathrm{C}) \xrightarrow {\mathrm{H} ^ {n} (u , g)} \mathrm{H} _ {\mathrm{A}} ^ {n} (u, \mathrm{C} ^ {\prime \prime}) \xrightarrow {\partial^ {n}} \mathrm{H} _ {\mathrm{A}} ^ {n + 1} (u, \mathrm{C} ^ {\prime}) \rightarrow \dots .
\]

Let \(\rho: A \to A'\) be a ring homomorphism, \(L'\) be the \(A'\) -module \(A' \otimes_A L\) , \(u': L' \to A'\) the linear form \(1 \otimes u\) . The canonical bijective homomorphism (III, p.~83, prop. 8)

\[
\psi : \Lambda_ {A ^ {\prime}} (A ^ {\prime} \otimes_ {A} L) \rightarrow A ^ {\prime} \otimes_ {A} \Lambda_ {A} (L)
\]

is an isomorphism of complexes of A'-modules. We deduce:

\begin{enumerate}
\def\labelenumi{\arabic{enumi})}
\tightlist
\item
  for every complex of A'-modules C', an isomorphism of complexes of A-modules
\end{enumerate}

\[
\mathbf {K} _ {\bullet} ^ {\mathrm{A} ^ {\prime}} \left(u ^ {\prime}, \mathrm{C} ^ {\prime}\right)\rightarrow \mathbf {K} _ {\bullet} ^ {\mathrm{A}} \left(u, \mathrm{C} ^ {\prime}\right),
\]

\begin{enumerate}
\def\labelenumi{(\arabic{enumi})}
\setcounter{enumi}{8}
\tightlist
\item
\end{enumerate}

composed of the diagram

\[
\mathrm{C} ^ {\prime} \otimes_ {\mathrm{A} ^ {\prime}} \left(\boldsymbol {\Lambda} _ {\mathrm{A} ^ {\prime}} \left(\mathrm{A} ^ {\prime} \otimes_ {\mathrm{A}} \mathrm{L}\right)\right) \xrightarrow {1 _ {\mathrm{C} ^ {\prime}} \otimes \psi} \mathrm{C} ^ {\prime} \otimes_ {\mathrm{A} ^ {\prime}} \mathrm{A} ^ {\prime} \otimes_ {\mathrm{A}} \boldsymbol {\Lambda} _ {\mathrm{A}} (\mathrm{L}) \xrightarrow {\varphi} \mathrm{C} ^ {\prime} \otimes_ {\mathrm{A}} \boldsymbol {\Lambda} _ {\mathrm{A}} (\mathrm{L})
\]

where φ is the canonical bijection (III, p.~85, prop. 14);

\begin{enumerate}
\def\labelenumi{\arabic{enumi})}
\setcounter{enumi}{1}
\tightlist
\item
  for every complex of A-modules C, an isomorphism of complexes of A'-modules
\end{enumerate}

\[
\mathbf {K} _ {\bullet} ^ {\mathrm{A}} (u, \mathrm{A} ^ {\prime} \otimes_ {\mathrm{A}} \mathrm{C}) \rightarrow \mathrm{A} ^ {\prime} \otimes_ {\mathrm{A}} \mathbf {K} _ {\bullet} ^ {\mathrm{A}} (u, \mathrm{C}),
\]

whence homomorphisms of A'-modules

\[
\mathrm{A} ^ {\prime} \otimes_ {\mathrm{A}} \mathrm{H} _ {n} ^ {\mathrm{A}} (u, \mathrm{C}) \rightarrow \mathrm{H} _ {n} ^ {\mathrm{A} ^ {\prime}} (u ^ {\prime}, \mathrm{A} ^ {\prime} \otimes_ {\mathrm{A}} \mathrm{C}),
\]

which are bijective when A' is flat over A (X, p.~66, cor. 2).

Let L' be an A-module, \(u': L' \to A\) a linear form, \(f: L \to L'\) an A-homomorphism such that \(u' \circ f = u\) . It follows from III, p.~161, formula (55), that the homomorphism \(\symbf{\Lambda}(f): \symbf{\Lambda}(\mathrm{L}) \to \symbf{\Lambda}(\mathrm{L}')\) satisfies \(d_u \circ \symbf{\Lambda}(f) = \symbf{\Lambda}(f) \circ d_{u'}\) , hence defines a morphism of complexes \(\symbf{\Lambda}(u): \mathbf{K}^{\mathbf{A}}(u) \to \mathbf{K}^{\mathbf{A}}(u')\) . If C is an A-complex, we deduce morphisms of complexes

\[
1 _ {C} \otimes \boldsymbol {\Lambda} (u): \mathbf {K} ^ {\mathrm{A}} (u, C) \rightarrow \mathbf {K} ^ {\mathrm{A}} \left(u ^ {\prime}, C\right) \quad \text { and }\quad \operatorname{Homgr} \left(\boldsymbol {\Lambda} (u), 1 _ {C}\right): \mathbf {K} _ {A} ^ {*} \left(u ^ {\prime}, C\right)\rightarrow \mathbf {K} _ {A} ^ {*} (u, C)
\]

If f is bijective, all these morphisms are isomorphisms.

\subsection{\texorpdfstring{Example 1: the complex \(\mathbf{S}(\mathbf{L})\otimes_{\mathbf{A}}\boldsymbol{\Lambda}(\mathbf{L})\)}{Example 1: the complex \textbackslash mathbf\{S\}(\textbackslash mathbf\{L\})\textbackslash otimes\_\{\textbackslash mathbf\{A\}\}\textbackslash boldsymbol\{\textbackslash Lambda\}(\textbackslash mathbf\{L\})}}

Let A be a ring, L an A-module, \(\mathbf{S}(\mathrm{L})\) its symmetric algebra, \(\mathbf{S}(\mathrm{L}) \otimes_{\mathrm{A}} \mathrm{L}\) be the \(\mathbf{S}(\mathrm{L})\) -module obtained by extension of scalars, \(u: \mathbf{S}(\mathrm{L}) \otimes_{\mathrm{A}} \mathrm{L} \to \mathbf{S}(\mathrm{L})\) the linear form such that \(u(s \otimes x) = sx\) for \(s \in \mathbf{S}(\mathrm{L})\) , \(x \in \mathrm{L}\) . By the canonical isomorphism of \(\mathbf{S}(\mathrm{L})\) -modules (III, p.~83, prop. 8)

\[
\boldsymbol {\Lambda} _ {\mathbf {S} (L)} (\mathbf {S} (L) \otimes_ {A} L) \rightarrow \mathbf {S} (L) \otimes_ {A} \boldsymbol {\Lambda} (L),
\]

the differential of the complex \(\mathbf{K}^{\mathbf{S}(\mathrm{L})}(u)\) is transported to the map

\[
d: \mathbf {S} (\mathrm{L}) \otimes_ {\mathrm{A}} \boldsymbol {\Lambda} (\mathrm{L}) \rightarrow \mathbf {S} (\mathrm{L}) \otimes_ {\mathrm{A}} \boldsymbol {\Lambda} (\mathrm{L})
\]

such that, for \(x_{1},\ldots,x_{p},y_{1},\ldots,y_{q}\) in L, we have

\[
\begin{array}{l} d \left(\left(x _ {1} \dots x _ {p}\right) \otimes \left(y _ {1} \wedge \dots \wedge y _ {q}\right)\right) \\ = \sum_ {i = 1} ^ {q} (- 1) ^ {i + 1} y _ {i} x _ {1} \dots x _ {p} \otimes \left(y _ {1} \wedge \dots \wedge y _ {i - 1} \wedge y _ {i + 1} \wedge \dots \wedge y _ {q}\right). \end{array}
\]

Note that \(d\) maps \(\mathbf{S}^{p}(\mathrm{L})\otimes\symbf{\Lambda}^{q}(\mathrm{L})\) to \(\mathbf{S}^{p+1}(\mathrm{L})\otimes\symbf{\Lambda}^{q-1}(\mathrm{L})\) , so that the complex of A-modules \(\mathbf{S}(\mathrm{L})\otimes\symbf{\Lambda}(\mathrm{L})\) decomposes as the direct sum of the complexes described by the following diagrams:

\[
\left(\mathcal {E} _ {n}\right): 0 \rightarrow \mathbf {S} ^ {0} L \otimes_ {\mathrm{A}} \boldsymbol {\Lambda} ^ {n} L \rightarrow \mathbf {S} ^ {1} L \otimes_ {\mathrm{A}} \boldsymbol {\Lambda} ^ {n - 1} L \rightarrow \dots \rightarrow \mathbf {S} ^ {n} L \otimes_ {\mathrm{A}} \boldsymbol {\Lambda} ^ {0} L \rightarrow 0, n \in \mathbf {N}.
\]

If the A-module L is the direct sum of a finite family \((\mathbf{L}_{i})_{i\in I}\) where I is totally ordered, the canonical bijection

\[
\bigotimes_ {i \in I} ^ {\varepsilon} \left(\mathbf {S} (L _ {i}) \otimes_ {A} \boldsymbol {\Lambda} (L _ {i})\right) \to \mathbf {S} (L) \otimes_ {A} \boldsymbol {\Lambda} (L)
\]

is an isomorphism of complexes of A-modules (this follows from prop. 2 of X, p.~148 or from formula (9) above).

PROPOSITION 3. --- If the A-module L is flat, then the sequences \((\mathcal{E}_n)\) above are exact for \(n > 0\) .

\begin{enumerate}
\def\labelenumi{\alph{enumi})}
\tightlist
\item
  Let us first note that, if \(p_{L}\) is the composite homomorphism
\end{enumerate}

\[
\mathbf {S} (\mathrm{L}) \otimes \boldsymbol {\Lambda} (\mathrm{L}) \xrightarrow {\alpha} \mathbf {S} ^ {0} (\mathrm{L}) \otimes \boldsymbol {\Lambda} ^ {0} (\mathrm{L}) \xrightarrow {\beta} \mathrm{A},
\]

where \(\alpha\) is the tensor product of the canonical projections and \(\beta\) the canonical isomorphism; it remains to prove that \(\mathbf{H}(p_{\mathrm{L}})\) is bijective.

\begin{enumerate}
\def\labelenumi{\alph{enumi})}
\setcounter{enumi}{1}
\item
  If L = 0 or if L = A, the proposition is evident.
\item
  Suppose L is free of finite rank; write it as a direct sum \(L_{1} \oplus \ldots \oplus L_{n}\) of free A-modules of rank 1. By the remark preceding the proposition, the complex \(\mathbf{S}(\mathrm{L}) \otimes \symbf{\Lambda}(\mathrm{L})\) is isomorphic to the tensor product of the n free complexes \(\mathbf{S}(\mathrm{L}_{i}) \otimes \symbf{\Lambda}(\mathrm{L}_{i})\) whose homology is free by b). By X, p.~79, cor. 4, the canonical homomorphism
\end{enumerate}

\[
\gamma : \bigoplus_ {i = 1} ^ {n} \mathrm{H} (\mathbf {S} (\mathrm{L} _ {i}) \otimes \boldsymbol {\Lambda} (\mathrm{L} _ {i})) \rightarrow \mathrm{H} (\mathbf {S} (\mathrm{L}) \otimes \boldsymbol {\Lambda} (\mathrm{L}))
\]

is bijective. By b) \(\mathrm{H}(p_{\mathrm{L}_i})\) is bijective for all \(i\) . Since \(\bigotimes_{i=1}^{n} \mathrm{H}(p_{\mathrm{L}_i}) = \mathrm{H}(p_{\mathrm{L}}) \circ \gamma\) ,

\[
\mathrm{H} (p _ {\mathrm{L}}) \text {   is   bijective }  .
\]

\begin{enumerate}
\def\labelenumi{\alph{enumi})}
\setcounter{enumi}{3}
\tightlist
\item
  In the general case, L is the inductive limit of a filtered inductive system \((\mathbf{L}_{i})_{i\in I}\) of free modules of finite rank (X, p.~14, th. 1). Since the canonical bijective homomorphism
\end{enumerate}

\[
\varinjlim \mathbf {S} (\mathrm{L} _ {i}) \otimes \boldsymbol {\Lambda} (\mathrm{L} _ {i}) \rightarrow \mathbf {S} (\mathrm{L}) \otimes \boldsymbol {\Lambda} (\mathrm{L})
\]

is an isomorphism of complexes, the proposition follows from X, p.~28, prop. 1.

Remarks. --- 1) We shall see below (X, p.~158, example) another proof of part c) above.

\begin{enumerate}
\def\labelenumi{\arabic{enumi})}
\setcounter{enumi}{1}
\tightlist
\item
  If A is a Q-algebra, the conclusion of Prop. 3 remains true without any hypothesis on L (cf. X, p. 206, Exercise 1).
\end{enumerate}

Let G be a group and \(\rho: G \to \mathbf{GL}(L)\) a linear representation of G in a flat A-module L. Then the \((\mathcal{E}_n)\) are exact sequences of linear representations. Suppose L is finitely generated projective, and denote by \(R_A(G)\) the ring of representations of G in finitely generated projective A-modules. It follows from Prop. 3 that we have in \(R_{A}(G)\) the relations

\[
\sum_ {i = 0} ^ {n} (- 1) ^ {i} \left[ \mathbf {S} ^ {i} (\mathrm{L}) \right] \left[ \boldsymbol {\Lambda} ^ {n - i} (\mathrm{L}) \right] = 0, \quad n > 0.\tag{10}
\]

If we consider the formal series

\[
s (\mathrm{T}) = \sum_ {i = 0} ^ {\infty} \left[ \mathbf {S} ^ {i} (\mathrm{L}) \right] \mathrm{T} ^ {i} \in \mathrm{R} _ {\mathbf {A}} (\mathrm{G}) [ [ \mathrm{T} ] ],
\]

\[
\lambda (\mathrm{T}) = \sum_ {i = 0} ^ {\infty} \left[ \boldsymbol {\Lambda} ^ {i} (\mathrm{L}) \right] \mathrm{T} ^ {i} \in \mathrm{R} _ {\mathbf {A}} (\mathrm{G}) [ [ \mathrm{T} ] ],
\]

relations (10) are written

\begin{enumerate}
\def\labelenumi{(\arabic{enumi})}
\setcounter{enumi}{10}
\tightlist
\item
\end{enumerate}

\[
s (\mathrm{T}) \lambda (- \mathrm{T}) = 1 _ {*}.
\]

\subsection{Example 2: the case of a free module}

Let \(k\) be a ring, \(\mathbf{M}\) a \(k\) -module, \(\mathbf{I}\) a set and \(p\) be an integer \(\geqslant 0\) . A map \(\boldsymbol{m}:1^{p}\to\mathbf{M}\) is said to be alternating if it satisfies the following two conditions:

\begin{enumerate}
\def\labelenumi{\alph{enumi})}
\tightlist
\item
  for every permutation \(\sigma \in \mathfrak{S}_p\) and every sequence \((\alpha_1, \ldots, \alpha_p) \in \mathrm{I}^p\) , we have
\end{enumerate}

\[
\boldsymbol {m} \left(\alpha_ {\sigma (1)}, \dots , \alpha_ {\sigma (p)}\right) = \varepsilon_ {\sigma} \boldsymbol {m} \left(\alpha_ {1}, \dots , \alpha_ {p}\right),
\]

\begin{enumerate}
\def\labelenumi{\alph{enumi})}
\setcounter{enumi}{1}
\tightlist
\item
  for every sequence \((\alpha_{1},\ldots,\alpha_{p})\in\mathbb{I}^{p}\) such that two of the indices \(\alpha_{1},\ldots,\alpha_{p}\) are equal, one has \(m(\alpha_{1},\ldots,\alpha_{p})=0\) .
\end{enumerate}

(In the case where I is a k-module and m is multilinear, one recovers the notion introduced in III, p.~80.)

Suppose I is finite and denote by \(C_{I}^{p}(M)\) the k-module of alternating maps from \(I^{p}\) to M.

Let \(L_{0}\) be a k-module, \((e_{i})_{i\in I}\) a family of elements of \(L_{0}\) ; we define two k-linear maps

\[
\begin{array}{l} g: \operatorname{Hom} _ {k} (\symbf {\Lambda} ^ {p} \mathrm{L} _ {0}, \mathrm{M}) \to \mathrm{C} _ {\mathrm{I}} ^ {p} (\mathrm{M}) \\ h: \mathrm{C} _ {\mathrm{I}} ^ {p} (\mathrm{M}) \to \mathrm{M} \otimes_ {k} \symbf {\Lambda} ^ {p} \mathrm{L} _ {0} \end{array}
\]

as follows: if \(f \in \operatorname{Hom}_k(\symbf{\Lambda}^p \mathbf{L}_0, \mathbf{M})\) , we set

\[
g (f) (\alpha_ {1}, \dots , \alpha_ {p}) = f (e _ {\alpha_ {1}} \wedge \dots \wedge e _ {\alpha_ {p}});
\]

let \(m \in C_{I}^{p}(M)\) , we define \(h(m) \in M \otimes_{k} \Lambda^{p} L_{0}\) . For every element \((\alpha_{1}, \ldots, \alpha_{p})\)

of \(I^{p}\) the element \(m(\alpha_{1},\ldots,\alpha_{p})\otimes(e_{\alpha_{1}}\wedge\ldots\wedge e_{\alpha_{p}})\) of \(\Lambda^{p}L_{0}\otimes_{k}M\) is zero if Card \(\{\alpha_{1},\ldots,\alpha_{p}\}<p\) and is independent of the order of the indices \(\alpha_{1},\ldots,\alpha_{p}\) if

\[
\text { Card } \{\alpha_ {1}, \dots , \alpha_ {p} \} = p.
\]

It depends only on the subset \(\mathbf{J} = (\alpha_{1},\ldots ,\alpha_{p})\) of I; denote it by \(h_\mathrm{J}(\pmb {m})\) ; we have \(h_\mathrm{J}(\pmb {m}) = 0\) if Card (J) \(< p\) ; we then set:

\[
h (\boldsymbol {m}) = \sum_ {\mathrm{J}} h _ {\mathrm{J}} (\boldsymbol {m}),
\]

where J ranges over the subsets of I having p elements.

Lemma 2. --- If the k-module \(\mathbf{L}_{0}\) is free with basis \((e_{i})_{i\in I}\) , the \(k\) -linear maps \(g\) and \(h\) are bijective.

This follows from III, p.~79, th. 1.

Now let M be a k-module, and \(\boldsymbol{x} = (x_{i})_{i \in I}\) a family of k-endomorphisms of M, pairwise commuting. Consider the polynomial ring \(\mathbf{A} = k[(\mathbf{X}_{i})_{i \in I}]\) and endow M with the A-module structure such that \(\mathrm{P}(\mathbf{X}_{i}) m = \mathrm{P}(x_{i}) m\) for \(P \in A\) and \(m \in M\) . Let, on the other hand, L be the free A-module \(A^{I}, (e_{i})_{i \in I}\) its canonical basis and \(u : L \to A\) the linear form which maps \(e_{i}\) over \(X_{i}\) for all \(i \in I\) . Consider the complexes of k-modules \(\mathbf{K}^{\mathbf{A}}(u, \mathbf{M})\) and \(\mathbf{K}_{\mathbf{A}}^{\bullet}(u, \mathbf{M})\) ; we have canonical isomorphisms

\[
\begin{array}{c} \mathbf {K} _ {\mathbf {A}} ^ {p} (u, \mathrm{M}) = \operatorname{Hom} _ {\mathbf {A}} \left(\boldsymbol {\Lambda} _ {\mathbf {A}} ^ {p} (\mathrm{A} ^ {\mathrm{I}}), \mathrm{M}\right) \to \operatorname{Hom} _ {k} \left(\boldsymbol {\Lambda} _ {k} ^ {p} (k ^ {\mathrm{I}}), \mathrm{M}\right), \\ \mathrm{M} \otimes_ {k} \boldsymbol {\Lambda} _ {k} ^ {p} (k ^ {\mathrm{I}}) \to \mathrm{M} \otimes_ {\mathbf {A}} \boldsymbol {\Lambda} _ {\mathbf {A}} ^ {p} (\mathrm{A} ^ {\mathrm{I}}) = \mathbf {K} _ {p} ^ {\mathbf {A}} (u, \mathrm{M}); \end{array}
\]

whence, by composition with the isomorphisms \(g\) and \(h\) , isomorphisms of \(k\) -modules

\[
\begin{array}{l} \theta^ {p}: \mathbf {K} _ {\mathrm{A}} ^ {p} (u, \mathrm{M}) \to \mathrm{C} _ {\mathrm{I}} ^ {p} (\mathrm{M}), \\ \theta_ {p}: \mathrm{C} _ {\mathrm{I}} ^ {p} (\mathrm{M}) \to \mathbf {K} _ {p} ^ {\mathrm{A}} (u, \mathrm{M}). \end{array}
\]

We denote by

\[
\begin{array}{l} \hat {c} ^ {p}: \mathbf {C} _ {\mathrm{I}} ^ {p} (\mathbf {M}) \to \mathbf {C} _ {\mathrm{I}} ^ {p + 1} (\mathbf {M}), \\ \hat {c} _ {p}: \mathbf {C} _ {\mathrm{I}} ^ {p} (\mathbf {M}) \to \mathbf {C} _ {\mathrm{I}} ^ {p - 1} (\mathbf {M}), \end{array}
\]

the \(k\) -homomorphisms obtained by transporting the differentials of \(\mathbf{K}_{\mathbf{A}}^{\bullet}(u, \mathbf{M})\) and \(\mathbf{K}_{\mathbf{A}}^{\mathbf{A}}(u, \mathbf{M})\) via the isomorphisms \(\theta\) . We have, for example:

\[
\left(\partial^ {p} \boldsymbol {m}\right) \left(\alpha_ {1}, \dots , \alpha_ {p + 1}\right) = \sum_ {j = 1} ^ {p + 1} (- 1) ^ {j + 1} x _ {\alpha_ {j}} \boldsymbol {m} \left(\alpha_ {1}, \dots , \alpha_ {j - 1}, \alpha_ {j + 1}, \dots , \alpha_ {p + 1}\right).\tag{12}
\]

The complex formed by the \(C_{I}^{p}(M)\) and the \(\partial^{p}\) (resp. the \(\partial_{p}\) ) is denoted \(\mathbf{K}^{\prime}(\mathbf{x},\mathbf{M})\) (resp. \(\mathbf{K}(\mathbf{x},\mathbf{M})\) ) and is called the ascending (resp. descending) Koszul complex

associated with the module M and the sequence of endomorphisms \((x_{1}, \ldots, x_{n})\) . We therefore have isomorphisms of complexes of k-modules

\[
\begin{array}{l} \theta^ {\bullet}: \mathbf {K} _ {\mathbf {A}} ^ {\bullet} (u, \mathbf {M}) \to \mathbf {K} ^ {\bullet} (x, \mathbf {M}), \\ \theta_ {\bullet}: \mathbf {K}. (x, \mathbf {M}) \to \mathbf {K} _ {\bullet} ^ {\mathbf {A}} (u, \mathbf {M}). \end{array}
\]

Remark. --- Conversely, let B be a \(k\) -algebra, L a free B-module with basis \((e_i)_{i \in I}\) , and M a B-module. The data of a linear form \(u: L \to B\) is equivalent to that of a family \(x = (x_i)_{i \in 1}\) of elements of B, via the relation \(x_i = u(e_i)\) . The complex of \(k\) -modules underlying \(\mathbf{K}_{\mathrm{B}}^{\bullet}(u, \mathrm{M})\) (resp. \(\mathbf{K}^{\mathrm{B}}(u, \mathrm{M})\) ) is then identified, via the isomorphism \(\theta^{\bullet}\) (resp. \(\theta\) .) with the Koszul complex \(\mathbf{K}'(x, \mathrm{M})\) (resp. \(\mathbf{K}^{\bullet}(x, \mathrm{M})\) ). For example \(\mathbf{K}^{\mathrm{B}}(u)\) is identified with \(\mathbf{K}'(x, \mathrm{B})\) .

We introduce, as in no. 1 (X, p.~147), the notations H.(x, M), H.(x, M), etc., and all the results of nos. 1 and 2 apply mutatis mutandis, the module A \(^{I}\) being free. For example, we have isomorphisms

\[
\begin{array}{l} \mathrm{H} _ {0} (\boldsymbol {x}, \mathrm{M}) \to \mathrm{M} / (\boldsymbol {x}) \mathrm{M} \\ \mathrm{H} ^ {0} (\boldsymbol {x}, \mathrm{M}) \to \mathrm{Hom} _ {\mathrm{A}} \left(\mathrm{A} / (\boldsymbol {x}), \mathrm{M}\right), \end{array}
\]

where \((\pmb{x})\) denotes the ideal \(\sum \mathbf{A}x_{i}\) of A. Also for example, if \(\mathbf{K}.(\pmb{x},\mathbf{A})\) is acyclic in degrees \(>0\) , we have isomorphisms

\[
\begin{array}{l}\mathrm{H} _ {r} (\boldsymbol {x}, \mathrm{M}) \rightarrow \operatorname{Tor} _ {r} ^ {\mathrm{A}} (k, \mathrm{M}),\\\operatorname{Ext} _ {\mathrm{A}} ^ {r} (k, \mathrm{M}) \rightarrow \mathrm{H} ^ {r} (\boldsymbol {x}, \mathrm{M}).\end{array}
\]

Finally, suppose I is (finite and) totally ordered, for example \(I = \{1, \ldots, n\}\) ; let us identify \(\symbf{\Lambda}^{n}(\mathbf{A}^{\mathrm{I}})\) to A by means of the basis element \(e_{1} \wedge \ldots \wedge e_{n}\) and let us translate the isomorphism \(\mathbf{K}_{p}^{\mathrm{A}}(u, \mathrm{M}) \to \mathbf{K}_{\mathrm{A}}^{n-p}(u, \mathrm{M})\) of X, p.~149. It becomes, by transport by \((\theta_{p})\) and \((\theta^{n-p})\) , the isomorphism

\[
\mathrm{C} _ {\mathrm{I}} ^ {p} (\mathbf {M}) \rightarrow \mathrm{C} _ {\mathrm{I}} ^ {n - p} (\mathbf {M})
\]

which associates with \(m \in C_{I}^{p}(M)\) the element \(\check{m}\) of \(C_{1}^{n-p}(M)\) such that

\[
\boldsymbol {m} \left(\alpha_ {1}, \dots , \alpha_ {p}\right) = \check {\boldsymbol {m}} \left(\beta_ {1}, \dots , \beta_ {n - p}\right)\tag{13}
\]

if \((\alpha_{1},\ldots,\alpha_{p},\beta_{1},\ldots,\beta_{n-p})\) is an even permutation of \(\{1,\ldots,n\}\) Let us also note that when \(I=\{1,\ldots,n\}\) , one can identify \(C_{I}^{p}(M)\) with the set of families \(m(\alpha_{1},\ldots,\alpha_{p})\) of elements of M where \(\alpha_{1}<\alpha_{2}<\ldots<\alpha_{p}\) ; formula (12) remains valid, as does relation (13).

Example. --- Take \(\mathbf{M} = k[\mathbf{T}_1, \ldots, \mathbf{T}_n]\) ; the Koszul complex \(\mathbf{K}\) ( \(\partial/\partial\mathbf{T}\) , \(\mathbf{M}\) ) associated with the sequence of endomorphisms ( \(\partial/\partial\mathbf{T}_1, \ldots, \partial/\partial\mathbf{T}_n\) ) is identified with the de Rham complex of \(k[x_{1},\ldots,x_{n}]\) over \(k\) (X, p.~44): to \(\boldsymbol{m}\in\mathbf{C}_{\{1,\ldots,n\}}^{p}(\mathbf{M})\) , one associates the differential form

\[
\omega (\boldsymbol {m}) = \sum_ {\alpha_ {1} <   \dots <   \alpha_ {p}} \boldsymbol {m} \left(\alpha_ {1}, \dots , \alpha_ {p}\right) d x _ {\alpha_ {1}} \wedge \dots \wedge d x _ {\alpha_ {p}},
\]

cf.~formula (12) and example 1, p.~44.

\subsection{Example 3: the case L = A}

If L = A, set \(u(1) = x \in \mathbf{A}\) . The complex \(\mathbf{K}(u)\) then has length 1, and we have \(\mathbf{K}_{0}(u) = \mathbf{K}_{1}(u) = \mathbf{A}\) and \(d_{1}(a) = xa\) , so for every A-module M, \(\mathbf{K}_{0}(u, \mathbf{M})\) , \(\mathbf{K}_{1}(u, \mathbf{M})\) , \(\mathbf{K}^{0}(u, \mathbf{M})\) and \(\mathbf{K}^{1}(u, \mathbf{M})\) are identified with M, the differentials

\[
d _ {1}: \mathbf {K} _ {1} (u, \mathrm{M}) \rightarrow \mathbf {K} _ {0} (u, \mathrm{M}) \quad \text { and }\quad d ^ {0}: \mathbf {K} ^ {0} (u, \mathrm{M}) \rightarrow \mathbf {K} ^ {1} (u, \mathrm{M})
\]

being \(m \mapsto xm\) . We therefore have isomorphisms

\[
\begin{array}{l} \mathrm{H} _ {0} (x, \mathrm{M}) \to \mathrm{A} / x \mathrm{A} \otimes_ {\mathrm{A}} \mathrm{M} \leftarrow \mathrm{H} ^ {1} (x, \mathrm{M}), \\ \mathrm{H} _ {1} (x, \mathrm{M}) \to \mathrm{Hom} _ {\mathrm{A}} (\mathrm{A} / x \mathrm{A}, \mathrm{M}) \leftarrow \mathrm{H} ^ {0} (x, \mathrm{M}). \end{array}
\]

Lemma 3. --- Let K be a complex such that \(\mathbf{K}_i = 0\) for \(i \neq 0, 1\) and let C be a complex and \(p\) an integer.

\begin{enumerate}
\def\labelenumi{\alph{enumi})}
\tightlist
\item
  If \(K\) is flat, we have for all \(p \in \mathbf{N}\) an exact sequence
\end{enumerate}

\[
0 \rightarrow \mathrm{H} _ {0} (\mathbf {K} \otimes_ {\mathbf {A}} \mathrm{H} _ {p} (\mathbf {C})) \rightarrow \mathrm{H} _ {p} (\mathbf {K} \otimes_ {\mathbf {A}} \mathbf {C}) \rightarrow \mathrm{H} _ {1} (\mathbf {K} \otimes_ {\mathbf {A}} \mathrm{H} _ {p - 1} (\mathbf {C})) \rightarrow 0.
\]

\begin{enumerate}
\def\labelenumi{\alph{enumi})}
\setcounter{enumi}{1}
\tightlist
\item
  If K is projective, we have for every \(p \in \mathbf{N}\) an exact sequence
\end{enumerate}

\[
\begin{array}{r l} 0 \to \mathrm{H} ^ {1} (\operatorname{Homgr} _ {\mathbf {A}} (\mathbf {K}, \mathrm{H} ^ {p - 1} (\mathbf {C}))) & \to \mathrm{H} ^ {p} (\operatorname{Homgr} _ {\mathbf {A}} (\mathbf {K}, \mathbf {C})) \\ & \to \mathrm{H} ^ {0} (\operatorname{Homgr} _ {\mathbf {A}} (\mathbf {K}, \mathrm{H} ^ {p} (\mathbf {C}))) \to 0. \end{array}
\]

Let us prove \(a\) , the proof of \(b\) ) being analogous. For every \(i\) , denote by \(\mathbf{K}_{(i)}\) the complex \(\mathbf{K}_i(-i)\) . We have an exact sequence of complexes, split as a sequence of A-modules

\[
0 \rightarrow \mathrm{K} _ {(0)} \xrightarrow {\alpha} \mathrm{K} \xrightarrow {\beta} \mathrm{K} _ {(1)} \rightarrow 0;
\]

the sequence

\[
0 \longrightarrow \mathrm{K} _ {(0)} \otimes_ {\mathrm{A}} \mathrm{C} \xrightarrow {\alpha \otimes 1} \mathrm{K} \otimes_ {\mathrm{A}} \mathrm{C} \xrightarrow {\beta \otimes 1} \mathrm{K} _ {(1)} \otimes_ {\mathrm{A}} \mathrm{C} \longrightarrow 0\tag{14}
\]

is exact, and since K is flat, the homomorphisms

\[
\gamma_ {0, p} (\mathbf {K} _ {(0)}, \mathbf {C}): \mathbf {K} _ {(0)} \otimes_ {\mathbf {A}} \mathrm{H} _ {p} (\mathbf {C}) \rightarrow \mathrm{H} _ {p} (\mathbf {K} _ {(0)} \otimes_ {\mathbf {A}} \mathbf {C})
\]

\[
\gamma_ {1, p - 1} (\mathbf {K} _ {(1)}, \mathbf {C}): \mathbf {K} _ {(1)} \otimes_ {\mathbf {A}} \mathrm{H} _ {p - 1} (\mathbf {C}) \rightarrow \mathrm{H} _ {p} (\mathbf {K} _ {(1)} \otimes_ {\mathbf {A}} \mathbf {C})
\]

are bijective (X, p.~66, cor. 2). Let us compute the connecting homomorphism \(\partial (\alpha \otimes 1, \beta \otimes 1)\) ; by definition, it maps the class of the cycle \(\sum a_i \otimes b_i\) to that of \(\sum da_i \otimes b_i\) , which means that

\[
\partial (\alpha \otimes 1, \beta \otimes 1) \circ \gamma (\mathbf {K} _ {(1)}, \mathbf {C}) = \gamma (\mathbf {K} _ {(0)}, \mathbf {C}) \circ (d _ {\mathbf {K}} \otimes 1).
\]

The exact homology sequence associated with (14) therefore takes the form

\[
\begin{array}{r l} \mathrm{K} _ {1} \otimes \mathrm{H} _ {p} (\mathrm{C}) & \xrightarrow {d _ {\mathbf {k}} \otimes 1} \mathrm{K} _ {0} \otimes \mathrm{H} _ {p} (\mathrm{C}) \longrightarrow \mathrm{H} _ {p} (\mathrm{K} \otimes \mathrm{C}) \\ & \longrightarrow \mathrm{K} _ {1} \otimes \mathrm{H} _ {p - 1} (\mathrm{C}) \xrightarrow {d _ {\mathbf {k}} \otimes 1} \mathrm{K} _ {0} \otimes \mathrm{H} _ {p - 1} (\mathrm{C}), \end{array}
\]

whence a).

Applying lemma 3, a) to the complex \(\mathbf{K}(u)\) , and using the commutation isomorphisms, we obtain:

PROPOSITION 4. --- For every complex C, we have exact sequences.

\[
0 \rightarrow \mathrm{A} / x \mathrm{A} \otimes_ {\mathrm{A}} \mathrm{H} _ {p} (\mathrm{C}) \rightarrow \mathrm{H} _ {p} (x, \mathrm{C}) \rightarrow \operatorname{Hom} _ {\mathrm{A}} (\mathrm{A} / x \mathrm{A}, \mathrm{H} _ {p - 1} (\mathrm{C})) \rightarrow 0.
\]

Corollary 1. --- For \(\mathbf{H}_{p}(x,\mathbf{C})=0\) it is necessary and sufficient that the homothety with ratio \(x\) to \(\mathbf{H}_{p}(\mathbf{C})\) be surjective and the homothety with ratio \(x\) to \(\mathbf{H}_{p-1}(\mathbf{C})\) be injective.

COROLLARY 2. --- Let \(x = (x_1, \ldots, x_n)\) be a sequence of elements of A, M an A-module, \(x'\) the sequence \((x_1, \ldots, x_{n-1})\) . We have exact sequences

\[
0 \rightarrow \mathrm{A} / x _ {n} \mathrm{A} \otimes_ {\mathrm{A}} \mathrm{H} _ {p} (x ^ {\prime}, \mathrm{M}) \rightarrow \mathrm{H} _ {p} (x, \mathrm{M}) \rightarrow \operatorname{Hom} _ {\mathrm{A}} (\mathrm{A} / x _ {n} \mathrm{A}, \mathrm{H} _ {p - 1} (x ^ {\prime}, \mathrm{M})) \rightarrow 0.
\]

Corollary 3. --- For \(\mathrm{H}_{i}(\mathbf{x},\mathbf{M})\) to be zero for all i \textgreater{} 0, it is necessary and sufficient that the homothety with ratio \(x_{n}\) to \(\mathrm{H}_{i}(\mathbf{x}^{\prime},\mathbf{M})\) be bijective for i \textgreater{} 0, and that the homothety with ratio \(x_{n}\) to \(\mathrm{M}/(\mathbf{x}^{\prime})\) M be injective.

\subsection{Completely intersecting families}

Let A be a ring, M an A-module, \(x = (x_{i})_{i \in I}\) a family of elements of A.

DEFINITION 2. --- The family x is said to be completely secant for M if one has \(\mathrm{H}_{i}(\mathbf{x},\mathbf{M})=0\) for i\textgreater0.

If I is finite, it amounts to the same thing (X, p.~150) to say that one has \(\mathbf{H}^i (\pmb {x},\mathbf{M}) = 0\) for \(i <   \mathrm{Card}\) (I).

The following proposition allows us to give examples of completely secant families.

PROPOSITION 5. — Let \(x = (x_1, \ldots, x_n)\) be a sequence of elements of A. If for \(i=1, \ldots, n\) the homothety of ratio \(x_i\) in the A-module \(M/(x_1 M + \cdots + x_{i-1} M)\) is injective, the sequence \(x\) is completely secant for M.

A sequence satisfying the conditions of the proposition is said to be regular for M, or M-regular. Note that this property generally depends on the order of the \(x_i\) . For example, the sequence (1, 0) is always M-regular, whereas the sequence (0, 1) is M-regular only if M is zero. On the other hand, whether a sequence is completely secant does not depend on the order of the terms.

Let us prove the proposition by induction on \(n\) , the case \(n=0\) being immediate. Set \(x'=(x_{1},...,x_{n-1})\) ; if the sequence \(x\) is M-regular, the sequence \(x'\) is M-regular and multiplication by \(x_{n}\) to \(\mathrm{M}/(\mathbf{x}')\mathrm{M}\) is injective; by the induction hypothesis, we have \(\mathrm{H}_{i}(\mathbf{x}',\mathrm{M})=0\) for \(i>0\) ; it then follows from cor. 3 of X, p.~157, that \(\mathrm{H}_{i}(\mathbf{x},\mathrm{M})=0\) for \(i>0\) .

Example. --- Let \(k\) be a ring; take \(\mathbf{A} = k[\mathbf{X}_1, \ldots, \mathbf{X}_n]\) and \(\boldsymbol{x} = (\mathbf{X}_1, \ldots, \mathbf{X}_n)\) . The sequence \(\boldsymbol{x}\) is \(\mathbf{A}\) -regular and the proposition gives back acyclicity in degrees \(>0\) of the complex \(\mathbf{S}_k(k^n) \otimes_k \boldsymbol{\Lambda}_k(k^n)\) (cf.~X, p.~152, prop. 3).

Similarly in the ring of formal power series \(\hat{\mathbf{A}} = k[[\mathbf{X}_1, \ldots, \mathbf{X}_n]]\) , the sequence \((\mathbf{X}_1, \ldots, \mathbf{X}_n)\) is \(\hat{\mathbf{A}}\) -regular, hence completely secant for \(\hat{\mathbf{A}}\) .

PROPOSITION 6. --- a) If \(\sum_{i\in I}x_i\mathbf{A} = \mathbf{A}\) , the family \((x_i)_{i\in I}\) is completely secant for M.

\begin{enumerate}
\def\labelenumi{\alph{enumi})}
\setcounter{enumi}{1}
\tightlist
\item
  Let \(x = (x_1, \ldots, x_n)\) be a sequence of elements of A. Let \((a_{ij}) \in \mathbf{GL}_n(\mathbf{A})\) ; set
\end{enumerate}

\[
y _ {i} = \sum_ {j} a _ {i j} x _ {j}.
\]

If the sequence x is completely secant for M, so is the sequence \((y_{1}, \ldots, y_{n})\) . c) Let \(k_{1}, \ldots, k_{n}\) be integers \(\geqslant 1\) ; for the sequence \((x_{1}^{k_{1}}, \ldots, x_{n}^{k_{n}})\) to be completely secant for M, it is necessary and sufficient that the sequence x be completely secant for M. Assertion a) follows from cor. 1, p.~148.

Let us prove \(b\) ). Let \(f: \mathbf{A}^n \to \mathbf{A}^n\) be the isomorphism defined by the matrix \(^t(a_{ij})\) ; it follows from X, p.~151, that \(1_{\mathbf{M}} \otimes \Lambda(f)\) is an isomorphism of the complex \(\mathbf{K}.(y, \mathbf{M})\) on the complex \(\mathbf{K}.(x, \mathbf{M})\) , whence \(b\) .

To prove \(c\) ), it suffices evidently to show that if \(k\) is an integer \(\geqslant 1\) , the sequence \((x_{1}, \ldots, x_{n-1}, x_{n}^{k})\) is completely secant for M if and only if the same holds for the sequence \(x\) . Let \(x' = (x_{1}, \ldots, x_{n-1})\) ; by cor. 3, p.~157, the first condition (resp. the second) means that the homothety of ratio \(x_{n}^{k}\) (resp. \(x_{n}\) ) is bijective in H \(_{i}(x', M)\) for \(i \geqslant 1\) and injective in M/(x') M. These two conditions are clearly equivalent, whence \(c\) ).

Remark. --- 1) The assertion analogous to c) for regular sequences is true (X, p.~207, exercise 5).

PROPOSITION 7. --- a) Let N be a flat A-module. If the family x is completely secant for M, it is so for M ⊗\_A N.

\begin{enumerate}
\def\labelenumi{\alph{enumi})}
\setcounter{enumi}{1}
\tightlist
\item
  Let \(0 \to \mathbf{M}' \to \mathbf{M} \to \mathbf{M}'' \to 0\) be an exact sequence of A-modules. If the family \(x\) is completely secant for \(\mathbf{M}'\) and for \(\mathbf{M}''\) , it is so for \(\mathbf{M}\) .
\end{enumerate}

The complex \(\mathbf{K}.\left(\mathbf{x},\mathbf{M}\otimes_{\mathbf{A}}\mathbf{N}\right)\) is isomorphic by definition to \(\mathbf{K}.\left(\mathbf{x},\mathbf{M}\right)\otimes_{\mathbf{A}}\mathbf{N}\) ; since N is flat, we deduce an isomorphism H. \((\mathbf{x},\mathbf{M})\otimes_{\mathbf{A}}\mathbf{N}\to\mathrm{H}.\left(\mathbf{x},\mathbf{M}\otimes_{\mathbf{A}}\mathbf{N}\right)\) (X, p.~66, cor. 2), whence \(a\) ).

The complex \(\mathbf{K}.\left(\mathbf{x}\right)\) being flat, we have an exact sequence of complexes

\[
0 \rightarrow \mathbf {K}. (x, \mathrm{M} ^ {\prime}) \rightarrow \mathbf {K}. (x, \mathrm{M}) \rightarrow \mathbf {K}. (x, \mathrm{M} ^ {\prime \prime}) \rightarrow 0;
\]

assertion b) follows from the associated exact homology sequence.

Remarks. --- 2) The assertions analogous to a) and b) for regular sequences are immediate.

\begin{enumerate}
\def\labelenumi{\arabic{enumi})}
\setcounter{enumi}{2}
\tightlist
\item
  If the family \(x\) is completely secant for \(A\) , the complex \(K.(x, A)\) defines a free resolution of the \(A\) -module \(A/x\) , with \(x = \sum_{i \in I} x_i A\) ; we therefore have for every integer \(i \geqslant 0\) and for every \(A\) -module \(M\) isomorphisms
\end{enumerate}

\[
\operatorname{Ext} _ {\mathbf {A}} ^ {n - i} (\mathbf {A} / \mathfrak {x}, \mathbf {M}) \rightarrow \mathrm{H} ^ {n - i} (\mathfrak {x}, \mathbf {M}) \rightarrow \mathrm{H} _ {i} (\mathfrak {x}, \mathbf {M}) \rightarrow \operatorname{Tor} _ {i} ^ {\mathbf {A}} (\mathbf {A} / \mathfrak {x}, \mathbf {M}).\tag{15}
\]

\begin{enumerate}
\def\labelenumi{\arabic{enumi})}
\setcounter{enumi}{3}
\tightlist
\item
  We say that the sequence \(\boldsymbol{x} = (x_1, \dots, x_n)\) is M-coregular if (denoting by \((x_i)_M\) the homothety with ratio \(x_i\) in M) the homothety of ratio \(x_i\) in the module
\end{enumerate}

\[
\operatorname{Ker} \left(x _ {1}\right) _ {\mathbf {M}} \cap \dots \cap \operatorname{Ker} \left(x _ {i - 1}\right) _ {\mathbf {M}}
\]

is surjective for \(i = 1, \ldots, n\) . We then have \(\mathbf{H}_{i}(\mathbf{x}, \mathbf{M}) = 0\) for \(i < n\) ; this is proved in the same manner as prop. 5.

\[
\mathrm{A} = k \left[ \mathrm{D} _ {1}, \dots , \mathrm{D} _ {n} \right],
\]

\[
\mathbf {Q} \text {-algebra},
\]

\(M = k[T_1, \ldots, T_n]\) , endowed with the structure of an A-module such that \(D_i P = \partial P / \partial T_i\) for all \(P \in M (1 \leqslant i \leqslant n)\) . One verifies immediately that the sequence \((D_1, \ldots, D_n)\) is M-coregular; taking into account the example on p.~155, one deduces that the de Rham complex of \(k[T_1, \ldots, T_n]\) over \(k\) is acyclic in degrees \textgreater{} 0.

\subsection{A criterion for completely secant sequences}

Let A be a ring, M an A-module, and x an ideal of A. The x-adic topology on M is the topology compatible with the group structure of M for which the set of submodules \(x^{r}\) M ( \(r \geqslant 0\) ) is a fundamental system of neighborhoods of zero (TG, III, p.~5, example). This topology is separated if and only if

\[
\bigcap_ {r \geqslant 0} x ^ {r} = 0.
\]

Suppose now that the ideal x is generated by a sequence \(\boldsymbol{x}=(x_{1},\ldots,x_{n})\)

of elements of A. Consider the graded A-module \(\bigoplus_{r\geqslant 0}x^{r}\) M and the graded A-homomorphism of degree 0

\[
a _ {\mathbf {M}} ^ {\mathbf {x}}: \mathrm{A} [ \mathrm{X} _ {1}, \dots , \mathrm{X} _ {n} ] \otimes_ {\mathbf {A}} \mathrm{M} \rightarrow \bigoplus_ {r \geqslant 0} \mathfrak {x} ^ {r} \mathrm{M}
\]

such that \(a_{\mathbf{M}}^{\star}(\mathbf{P} \otimes m) = \mathbf{P}(x_1, \ldots, x_n) \, m\) if \(\mathbf{P}\) is a homogeneous polynomial in \(\mathbf{X}_1, \ldots, \mathbf{X}_n\) and \(m \in \mathbf{M}\) . Let \(\mathfrak{d}\) be the ideal of \(\mathbf{A}[\mathbf{X}_1, \ldots, \mathbf{X}_n]\) generated by the elements \((x_i \, \mathbf{X}_j - x_j \, \mathbf{X}_i)\) for \(1 \leqslant i < j \leqslant n\) . We have \(\mathbf{P}(x_1, \ldots, x_n) = 0\) if \(\mathbf{P} \in \mathfrak{d}\) , so that \(a_{\mathbf{M}}^{\star}\) gives, by passing to the quotient, a graded A-homomorphism of degree 0

\[
\alpha_ {\mathrm{M}} ^ {x}: \left(\mathrm{A} \left[ \mathrm{X} _ {1}, \dots , \mathrm{X} _ {n} \right] / \mathfrak {d}\right) \otimes_ {\mathrm{A}} \mathrm{M} \rightarrow \bigoplus_ {r \geqslant 0} x ^ {r} \mathrm{M}.
\]

By taking the tensor product with A/x, we deduce from \(\alpha_{\mathrm{M}}^{\mathbf{x}}\) a graded A-homomorphism of degree 0

\[
\beta_ {\mathrm{M}} ^ {x}: (\mathrm{A} / \mathfrak {x}) [ \mathrm{X} _ {1}, \dots , \mathrm{X} _ {n} ] \otimes_ {\mathrm{A}} \mathrm{M} \rightarrow \bigoplus_ {r \geqslant 0} (\mathfrak {x} ^ {r} \mathrm{M} / \mathfrak {x} ^ {r + 1} \mathrm{M}).
\]

Homomorphisms \(a_{\mathbf{M}}^{x}\) , \(\alpha_{\mathbf{M}}^{x}\) and \(\beta_{\mathbf{M}}^{x}\) are surjective.

THEOREM 1. --- Consider the following conditions:

\begin{enumerate}
\def\labelenumi{(\roman{enumi})}
\item
  The sequence x is M-regular (X, p.~158).
\item
  The sequence x is completely secant for M (X, p.~157, def. 2).
\item
  One has \(\mathbf{H}_1(x, \mathbf{M}) = 0\) .
\item
  \(L^{\prime}\) homomorphism \(\alpha_{\mathbf{M}}^{x}:(\mathbf{A}[\mathbf{X}_{1},\dots,\mathbf{X}_{n}]/\mathfrak{d})\otimes_{\mathbf{A}}\mathbf{M}\to\bigoplus x^{r}\mathbf{M}\) is bijective.
\item
  The homomorphism \(\beta_{\mathbf{M}}^{\mathbf{x}}:(\mathbf{A}/\mathfrak{x})[\mathbf{X}_{1},..., \mathbf{X}_{n}]\otimes_{\mathbf{A}}\mathbf{M}\to\bigoplus_{r\geqslant0}(\mathfrak{x}^{r}\mathbf{M}/\mathfrak{x}^{r+1}\mathbf{M})\) is bijective.
\end{enumerate}

We then have the implications (i)⇒(ii)⇒(iii)⇔(iv)⇒(v). If for \(1 \leqslant i \leqslant n\) the A-module \(\mathrm{M}/(x_{1}\mathrm{M} + \cdots + x_{i-1}\mathrm{M})\) is separated for the x-adic topology, conditions (i) through (v) are equivalent.

The theorem will be proved in nos. 8 and 9.

COROLLARY 1. --- If A is Noetherian, if the A-module M is finitely generated, and if the \(x_{i}\) belong to the radical of A, conditions (i) through (v) of the theorem are equivalent. Indeed, on each of the modules \(\mathbf{M}/(x_{1} \mathbf{M} + \cdots + x_{i-1} \mathbf{M})\) the x-adic topology is separated (AC III, § 3, no. 3, prop. 6).

COROLLARY 2. --- Suppose that A is a positively graded ring, M a graded A-module bounded below, and the \(x_i\) homogeneous elements of degree \textgreater{} 0 of A. Then conditions (i) through (v) of the theorem are equivalent.

The x-adic topology is indeed separated for any graded A-module N that is bounded below, since if \(N_{n}=0\) for \(n<n_{0}\) we have \(x^{a}N\subset\sum_{j\geqslant n_{0}+a}N_{j}\) for all \(a\geqslant0\) .

COROLLARY 3. --- Suppose that the modules \(\mathbf{M}/(x_1\mathbf{M} + \cdots + x_{i-1}\mathbf{M})\) are separated for the \(\mathbf{x}\) -adic topology ( \(1 \leqslant i \leqslant n\) ); let \(p\) be an integer between 1 and \(n\) . For the sequence \(\mathbf{x}\) to be completely secant for \(\mathbf{M}\) , it is necessary and sufficient that the sequence \((x_1, \ldots, x_p)\) be completely secant for \(\mathbf{M}\) and that the sequence \((x_{p+1}, \ldots, x_n)\) be completely secant for \(\mathbf{M}/(x_1\mathbf{M} + \cdots + x_p\mathbf{M})\) .

Indeed, the corollary is obvious if in the statement one replaces "completely secant sequences" by "regular sequences"; and the two notions coincide here by the theorem.

Remark. --- Let \(x = (x_1, \ldots, x_n)\) be a sequence of elements of A such that \(H_1(x, A) = 0\) ; then the kernel of the surjective homomorphism \(u: A^n \to x\) such that \(u(\sum a_i e_i) = \sum a_i x_i\) is generated by the elements \(X_j e_i - X_i e_j\) ; consequently, the A-algebra \(A[X_1, \ldots, X_n] / \mathfrak{d}\) is isomorphic to the symmetric algebra \(S_A(x)\) (III, p.~69, prop. 4). It therefore follows from Theorem 1 that the algebra homomorphism \(S_A(x) \to \bigoplus_n x^n\) deduced from the canonical injection of \(x\) to \(\bigoplus_n x^n\) is an isomorphism. The same holds for the homomorphism \(S_A(x / x^2) \to \bigoplus_n x^n / x^{n+1}\) .

\subsection{Proof of Theorem 1: first part}

The implication (i) \(\Rightarrow\) (ii) has already been proved (X, p.~157, prop. 5). The implication (ii) \(\Rightarrow\) (iii) is evident; likewise for (iv) \(\Rightarrow\) (v), since \(\beta_{\mathbf{M}}^{x}\) is identified with \(\alpha_{\mathbf{M}}^{x} \otimes 1_{\mathbf{A}/\mathfrak{x}}\) .

To show that (iv) implies (iii), consider the homomorphism \((\alpha_{M}^{x})_{1}\) induced on the degree-1 components. Let E denote the graded A-module \(A[X_{1},\ldots,X_{n}]\) ; the A-module \(E_{1}\) is free with basis \(X_{1},\ldots,X_{n}\) and \(\mathfrak{d}_{1}\) is the sub-A-module of \(E_{1}\) generated by the elements \((x_{i}X_{j}-x_{j}X_{i})\) for \(1\leqslant i<j\leqslant n\) . Consequently, \(((E/\mathfrak{d})\otimes_{A}M)_{1}\) is identified with \(K_{1}(x,M)/B_{1}(K.(x,M))\) , the homomorphism \((\alpha_{M}^{x})_{1}\) being identified with the map from \(K_{1}(x,M)/B_{1}(K.(x,M))\) over \(B_{0}(K.(x,M))\) induced by \(d_{1}\) . Thus the vanishing of \(H_{1}(x,M)\) is equivalent to the injectivity of \((\alpha_{M}^{x})_{1}\) , whence the implication (iv) \(\Rightarrow\) (iii).

To prove that (iii) implies (iv), we shall use the following lemma:

Lemma 4. --- Let \(\mathbf{A}_{0}\) be the ring \(\mathbf{Z}[\mathrm{T}_{1},\ldots,\mathrm{T}_{n}]\) , and let \(u:\mathbf{A}_{0}[\mathrm{X}_{1},\ldots,\mathrm{X}_{n}] \to \mathbf{A}_{0}[\mathrm{U}]\) be the homomorphism of \(\mathbf{A}_{0}\) -algebras such that \(u(\mathbf{X}_{i}) = \mathbf{T}_{i}\) U. The kernel of \(u\) is the ideal \(\mathfrak{d}_{0}\) of \(\mathbf{A}_{0}[\mathbf{X}_{1},\ldots,\mathbf{X}_{n}]\) generated by the elements ( \(\mathrm{T}_{i}\mathbf{X}_{j}-\mathrm{T}_{j}\mathbf{X}_{i}\) ) for \(1 \leqslant i < j \leqslant n\) . If t is the ideal of \(\mathbf{A}_{0}\) generated by ( \(\mathrm{T}_{1},\ldots,\mathrm{T}_{n}\) ), \(u\) induces an isomorphism

\[
\overline {{u}}: \mathrm{A} _ {0} [ \mathrm{X} _ {1}, \dots , \mathrm{X} _ {n} ] / \mathfrak {d} _ {0} \rightarrow \bigoplus_ {r \geqslant 0} t ^ {r}.
\]

It obviously suffices to prove the first assertion. For every sequence of natural integers \(\alpha = (\alpha_{1},\ldots,\alpha_{n})\) and every integer \(k\in[0,n]\) , denote by \(\mathbf{P}_{\alpha,k}\) the monomial

\[
\mathrm{T} _ {1} ^ {\alpha_ {1}} \dots \mathrm{T} _ {k} ^ {\alpha_ {k}} \mathrm{X} _ {k + 1} ^ {\alpha_ {k + 1}} \dots \mathrm{X} _ {n} ^ {\alpha_ {n}};
\]

let N be the sub-Z-module of \(A_{0}[X_{1}, \ldots, X_{n}]\) generated by the \(P_{\alpha,k}\) for \(\alpha \in N^{n}\) and \(0 \leqslant k \leqslant n\) . We shall show that the restriction of u to N is injective and that \(A_{0}[X_{1}, \ldots, X_{n}] = \mathfrak{d}_{0} + N\) ; since one has \(d_{0} \subset Ker u\) , this will imply the lemma.

Observe that N is generated by the set S consisting of \(P_{0,0} = 1\) and of those \(P_{\alpha,k}\) for which \(\alpha_k \neq 0\) . To prove the injectivity of the restriction of u to N, it suffices to show that two distinct elements of S have as images under u distinct monomials in \(A_0[U]\) . Now one has \(u(P_{\alpha,k}) = T^\alpha U^{\sum_{i \geq k} \alpha_i}\) , so that the equality \(u(P_{\alpha,k}) = u(P_{\alpha',k'})\) implies \(\alpha = \alpha'\) and \(\sum_{i \geq k} \alpha_i = \sum_{i \geq k'} \alpha_i\) . Suppose that \(P_{\alpha,k}\) and \(P_{\alpha',k'}\) belong to S. If \(\alpha = 0\) , one then has \(k = k' = 0\) ; if \(\alpha \neq 0\) , we obtain \(k = k'\) since \(\alpha_k\) and \(\alpha_{k'}\) are nonzero, hence the result.

Let us show that every monomial \(\mathbf{T}^{\alpha}\mathbf{X}^{\beta}\in\mathbf{A}_{0}[\mathbf{X}_{1},..., \mathbf{X}_{n}]\) is congruent modulo \(\mathfrak{d}_{0}\) to a \(\mathbf{P}_{\eta,k}\) . For every sequence \(\lambda\in\mathbf{N}^{n}\) , we denote by \(i(\lambda)\) (resp. \(j(\lambda)\) ) the smallest (resp. the largest) integer \(k\in[1,n]\) such that \(\lambda_{k}\neq0\) . In the set of monomials \(\mathbf{T}^{\gamma}\mathbf{X}^{\delta}\) that are congruent to \(\mathbf{T}^{\alpha}\mathbf{X}^{\beta}\) modulo \(\mathfrak{d}_{0}\) let us choose one for which the rational integer \(j(\gamma)-i(\delta)\) is minimal; let us show that we then have \(j(\gamma)-i(\delta)<0\) . Suppose that we have \(j(\gamma)\geqslant i(\delta)\) ; denote by \(j=j(\gamma)\) , \(i=i(\delta)\) , and let \(\varepsilon=\inf(\gamma_{j},\delta_{i})\) . Since \((\mathbf{T}_{j}^{\varepsilon}\mathbf{X}_{i}^{\varepsilon}-\mathbf{T}_{i}^{\varepsilon}\mathbf{X}_{j}^{\varepsilon})\) is divisible by \((\mathbf{T}_{j}\mathbf{X}_{i}-\mathbf{T}_{i}\mathbf{X}_{j})\) , hence belongs to \(\mathfrak{d}_{0}\) , we see that \(\mathbf{T}^{\gamma}\mathbf{X}^{\delta}\) is congruent modulo \(\mathfrak{d}_{0}\) to \(\mathbf{T}^{\gamma'}\mathbf{X}^{\delta'}\) , where \(\gamma_{i}'=\gamma_{i}+\varepsilon\) , \(\gamma_{j}'=\gamma_{j}-\varepsilon\) , \(\gamma_{k}'=\gamma_{k}\) for \(k\neq i,j\) , and \(\delta_{i}'=\delta_{i}-\varepsilon\) , \(\delta_{j}'=\delta_{j}+\varepsilon\) , \(\delta_{k}'=\delta_{k}\) for \(k\neq i,j\) . Since \(\gamma_{j}'\) or \(\delta_{i}'\) is zero, we have \(j(\gamma')-i(\delta')<j(\gamma)-i(\delta)\) which contradicts the minimal character of \(j(\gamma)-i(\delta)\) .

Consequently we have \(j(\gamma) < i(\delta)\) , whence \(\mathbf{T}^{\gamma}\mathbf{X}^{\delta} \in \mathbf{N}\) , which completes the proof of the lemma.

Let us prove that (iii) implies (iv). Consider the ring \(\mathbf{A}_0 = \mathbf{Z}[\mathrm{T}_1, \ldots, \mathrm{T}_n]\) and the ideal t of \(\mathbf{A}_0\) generated by \(\mathrm{T}_1, \ldots, \mathrm{T}_n\) . Endow M with the structure of \(\mathbf{A}_0\) -module for which \(\mathrm{T}_i m = x_i m\) for \(m \in \mathbf{M}\) , \(1 \leqslant i \leqslant n\) . By X, p.~155, \(\mathrm{H}_i(x, \mathbf{M})\) is canonically identified with \(\mathbf{H}_i(\mathbf{T}, \mathbf{M})\) .

Since the sequence \(\mathbf{T}\) is regular for \(\mathbf{A}_{0}\) it follows from implication (i) \(\Rightarrow\) (ii) that the complex \(\mathbf{K}.\left(\mathbf{T},\mathbf{A}_{0}\right)\) defines a free resolution of the \(\mathbf{A}_{0}\) -module \(\mathbf{A}_{0}/t\) . Note that this one identifies with \(\mathbf{Z}\) , equipped with the structure of \(\mathbf{A}_{0}\) -module such that \(\mathrm{T}_{i}\mathbf{Z}=0\) for \(1\leqslant i\leqslant n\) . Thus condition (iii) is equivalent to \(\operatorname{Tor}_{1}^{\mathbf{A}_{0}}(\mathbf{M},\mathbf{Z})=0\) .

Let us show that (iii) implies \(\mathrm{Tor}_{1}^{\mathbf{A}_{0}}(\mathbf{M},\mathbf{A}_{0}/t^{r})=0\) for all \(r\geqslant1\) . This follows from the preceding for \(r=1\) . In the general case, consider the exact sequence

\[
0 \rightarrow t ^ {r} / t ^ {r + 1} \rightarrow A _ {0} / t ^ {r + 1} \rightarrow A _ {0} / t ^ {r} \rightarrow 0.\tag{16}
\]

The \(A_{0}\) -module \(t^{r}/t^{r+1}\) is isomorphic to a finite product of copies of Z; we therefore have \(\operatorname{Tor}_{1}^{\mathbf{A}_{0}}(\mathbf{M}, t^{r}/t^{r+1}) = 0\) . From this, by induction on r, we deduce that

\[
\operatorname{Tor} _ {1} ^ {\mathrm{A} _ {0}} \left(\mathrm{M}, \mathrm{A} _ {0} / t ^ {r}\right) = 0 \quad \text { for   all } r.
\]

The exact sequence (16) therefore yields, by tensor product with M, an exact sequence

\[
0 \rightarrow (t ^ {r} / t ^ {r + 1}) \otimes_ {\mathrm{A} _ {0}} \mathrm{M} \rightarrow \mathrm{M} / x ^ {r + 1} \mathrm{M} \rightarrow \mathrm{M} / x ^ {r} \mathrm{M} \rightarrow 0\tag{17}
\]

whence an isomorphism of \((\mathfrak{t}^r /\mathfrak{t}^{r + 1})\otimes_{\mathbf{A}_0}\mathbf{M}\) over \(\mathfrak{x}^r\mathbf{M} / \mathfrak{x}^{r + 1}\mathbf{M}\) . Considering the exact sequence \(0\to \mathfrak{t}^{r + 1}\to \mathfrak{t}^r\to \mathfrak{t}^r /\mathfrak{t}^{r + 1}\to 0\) , one then sees by induction on \(r\) that the map \(m_r:\mathfrak{t}^r\otimes_{\mathbf{A}_0}\mathbf{M}\rightarrow \mathfrak{x}^r\mathbf{M}\) , induced by the action of \(\mathbf{A}_0\) to \(\mathbf{M}\) is an isomorphism.

To prove (iv), it remains to observe that the diagram

\[
\begin{array}{c} \left(\mathrm{A} _ {0} \left[ \mathrm{X} _ {1}, \dots , \mathrm{X} _ {n} \right] / \mathfrak {d} _ {0}\right) \otimes_ {\mathrm{A} _ {0}} \mathrm{M} \xrightarrow {\bar {u} \otimes 1 _ {\mathrm{M}}} \bigoplus_ {r \geqslant 0} \left(\mathrm{t} ^ {r} \otimes_ {\mathrm{A} _ {0}} \mathrm{M}\right) \\ \Bigg \downarrow^ {e} \\ ^ {\prime} (\mathrm{A} [ \mathrm{X} _ {1}, \dots , \mathrm{X} _ {n} ] / \mathfrak {d}) \otimes_ {\mathrm{A}} \mathrm{M} \xrightarrow {\alpha_ {\mathrm{M}} ^ {*}} \bigoplus_ {r \geqslant 0} x ^ {r} \mathrm{M} \end{array}
\]

where e is the canonical isomorphism of extension of scalars (II, p.~85, prop. 6), is commutative, and to apply Lemma 4.

\subsection{Proof of Theorem 1: second part}

Consider again the exact sequence

\[
0 \longrightarrow t ^ {r} / t ^ {r + 1} \xrightarrow {i _ {r}} A _ {0} / t ^ {r + 1} \xrightarrow {p _ {r}} A _ {0} / t ^ {r} \longrightarrow 0\tag{16}
\]

and the associated exact sequence of torsion modules

\[
\begin{array}{l} \text {Tor} _ {1} ^ {\mathrm{A} _ {0}} (\mathrm{A} _ {0} / t ^ {r + 1}, \mathrm{M}) \longrightarrow \mathrm{Tor} _ {1} ^ {\mathrm{A} _ {0}} (\mathrm{A} _ {0} / t ^ {r}, \mathrm{M}) \\ \longrightarrow (t ^ {r} / t ^ {r + 1}) \otimes_ {\mathrm{A} _ {0}} \mathrm{M} \xrightarrow {i , \otimes 1 _ {\mathrm{M}}} (\mathrm{A} _ {0} / t ^ {r + 1}) \otimes_ {\mathrm{A} _ {0}} \mathrm{M} \xrightarrow {p , \otimes 1 _ {\mathrm{M}}} (\mathrm{A} _ {0} / t ^ {r}) \otimes_ {\mathrm{A} _ {0}} \mathrm{M} \longrightarrow 0. \end{array} \tag {18}
\]

The kernel of \(p_r \otimes 1_M\) is identified with \(x^r M/x^{r+1} M\) ; moreover, \(t^r/t^{r+1}\) is annihilated by the \(T_i\) and is identified with the homogeneous component of degree \(r\) of \(A_0\) , so that the homomorphism \((t^r/t^{r+1}) \otimes_{A_0} M \to x^r M/x^{r+1} M\) deduced from \(i_r \otimes 1_M\) is identified with the homogeneous component of degree \(r\) of the homomorphism \(\beta_M^x\) . It therefore follows from the exact sequence (18) that condition (v) is equivalent to

(v'): the homomorphism \(\operatorname{Tor}_{1}^{\mathbf{A}_{0}}(p_{r},1_{\mathbf{M}})\) : \(\operatorname{Tor}_{1}^{\mathbf{A}_{0}}(\mathbf{A}_{0}/t^{r+1},\mathbf{M})\to\operatorname{Tor}_{1}^{\mathbf{A}_{0}}(\mathbf{A}_{0}/t^{r},\mathbf{M})\) is surjective for all \(r\geqslant1\) .

It remains for us to prove the implication (v) \(\Rightarrow\) (i) when the modules

\[
\mathbf {M} / (x _ {1} \mathbf {M} + \dots + x _ {i - 1} \mathbf {M})
\]

are separated for the \(\mathfrak{x}\) -adic topology ( \(1 \leqslant i \leqslant n\) ). Let \(\overline{\mathbf{M}}\) be the A-module \(\mathbf{M}/x_1 \mathbf{M}\) . By definition, the sequence \(\mathbf{x}\) is regular for \(\mathbf{M}\) if and only if \((x_1)_\mathbf{M}\) is injective and the sequence \(\mathbf{x}' = (x_2, \ldots, x_n)\) is regular for \(\overline{\mathbf{M}}\) . Reasoning by induction on \(n\) , it therefore suffices to prove that, if \(\mathbf{M}\) is separated for the \(\mathfrak{x}\) -adic topology and if \(\beta_{\mathbf{M}}^{\mathbf{x}}\) is bijective, then \((x_1)_\mathbf{M}\) is injective and \(\beta_{\mathbf{M}}^{\mathbf{x}'}\) is bijective. Now the bijectivity of \(\beta_{\mathbf{M}}^{\mathbf{x}}\) implies in particular that the homothety with ratio \(x_1\) to \(\bigoplus_{r} x^r \mathbf{M}/x^{r+1} \mathbf{M}\) is injective, hence that \(\text{Ker } (x_1)_\mathbf{M} \subset \bigcap_{i} x^i \mathbf{M}\) and consequently that \((x_1)_\mathbf{M}\) is injective if the \(\mathfrak{x}\) -adic topology on \(\mathbf{M}\) is separated.

We are thus reduced to proving that if \((x_{1})_{\mathbf{M}}\) is injective and if M satisfies the condition \((\mathbf{v}')\) , then \(\overline{M}\) satisfies the condition \((\mathbf{v}')\) relative to the sequence \(x'\) .

By hypothesis, we have an exact sequence

\[
0 \longrightarrow \mathrm{M} \xrightarrow {(x _ {1}) _ {M}} \mathrm{M} \longrightarrow \overline {{\mathrm{M}}} \longrightarrow 0;
\]

let us set \(\overline{\mathbf{A}}_{0} = \mathbf{A}_{0}/\mathbf{T}_{1}\mathbf{A}_{0},\quad\overline{\mathbf{t}} = \mathbf{t}\overline{\mathbf{A}}_{0}\) . Let \(q:L\to M\) be a free resolution of the \(\mathbf{A}_{0}\) -module M; since the homothety of ratio \(\mathbf{T}_{1}\) is injective in \(\mathbf{A}_{0}\) , it is injective in L, and we have an exact sequence of complexes

\[
0 \longrightarrow \mathrm{L} \xrightarrow {(x _ {1}) _ {\mathrm{L}}} \mathrm{L} \longrightarrow \overline {{\mathrm{L}}} \longrightarrow 0
\]

with \(\overline{\mathbf{L}} = \mathbf{L} / x_1\mathbf{L}\) , and a commutative diagram

\[
\begin{array}{c} 0 \longrightarrow \mathrm{L} \stackrel {(x _ {1}) _ {\mathrm{L}}} {\longrightarrow} \mathrm{L} \longrightarrow \overline {{\mathrm{L}}} \longrightarrow 0 \\ \Big \downarrow^ {q} \quad \Big \downarrow^ {q} \quad \Big \downarrow^ {\overline {{q}}} \\ 0 \longrightarrow \mathrm{M} \stackrel {(x _ {1}) _ {\mathrm{M}}} {\longrightarrow} \mathrm{M} \longrightarrow \overline {{\mathrm{M}}} \longrightarrow 0. \end{array}
\]

Since \(q\) is a homologism, \(\overline{q}:\overline{\mathbf{L}}\to\overline{\mathbf{M}}\) is a free resolution of the \(\overline{\mathbf{A}}_{0}\) -module \(\overline{\mathbf{M}}\) (cor. 1, p.~30). For every \(\overline{\mathbf{A}}_{0}\) -module P, there is a canonical isomorphism

\[
\mathbf {P} \otimes_ {\mathbf {A} _ {0}} \mathbf {L} \rightarrow \mathbf {P} \otimes_ {\overline {{{\mathbf {A}}}} _ {0}} \overline {{{\mathbf {L}}}},
\]

whence, by passing to homology, an isomorphism

\[
\varphi_ {P}: \operatorname{Tor} _ {1} ^ {\mathbf {A} _ {0}} (\mathbf {P}, \mathbf {M}) \rightarrow \operatorname{Tor} _ {1} ^ {\overline {{\mathbf {A}}} _ {0}} (\mathbf {P}, \overline {{\mathbf {M}}}).
\]

If \(u: \mathbf{P} \to \mathbf{P}'\) is a \(\overline{\mathbf{A}}\) -homomorphism, then

\[
\varphi_ {\mathrm{P} ^ {\prime}} \circ \operatorname{Tor} _ {1} ^ {\mathbf {A} _ {0}} (u, 1 _ {\mathrm{M}}) = \operatorname{Tor} _ {1} ^ {\overline {{\mathbf {A}}} _ {0}} (u, 1 _ {\mathrm{M}}) \circ \varphi_ {\mathrm{P}}.
\]

That being so, suppose condition (v') is satisfied for M, and let us prove that it holds for \(\overline{M}\) . Let \(r\) be an integer \(\geqslant 1\) ; we have a commutative diagram with exact rows

\[
\begin{array}{c} 0 \longrightarrow \mathrm{A} _ {0} / t ^ {r} \xrightarrow {\mathrm{T} _ {1}} \mathrm{A} _ {0} / t ^ {r + 1} \longrightarrow \overline {{\mathrm{A}}} _ {0} / \overline {{t}} ^ {r + 1} \longrightarrow 0 \\ \downarrow^ {p _ {r - 1}} \quad \downarrow^ {p _ {r}} \quad \downarrow^ {\bar {p} _ {r}} \\ 0 \longrightarrow \mathrm{A} _ {0} / t ^ {r - 1} \xrightarrow {\mathrm{T} _ {1}} \mathrm{A} _ {0} / t ^ {r} \longrightarrow \overline {{\mathrm{A}}} _ {0} / \overline {{t}} ^ {r} \longrightarrow 0 \end{array}
\]

from which we deduce a commutative diagram with exact rows

\[
\begin{array}{c c c c c} \operatorname{Tor} _ {1} ^ {\mathbf {A} _ {0}} (\mathbf {A} _ {0} / t ^ {r + 1}, \mathbf {M}) & \longrightarrow & \operatorname{Tor} _ {1} ^ {\mathbf {A} _ {0}} (\overline {{\mathbf {A}}} _ {0} / \overline {{t}} ^ {r + 1}, \mathbf {M}) & \longrightarrow & \mathbf {M} / x ^ {r}   \mathbf {M} \xrightarrow {x _ {1}} \mathbf {M} / x ^ {r + 1}   \mathbf {M} \\ \operatorname{Tor} _ {1} (p _ {r}, 1) & \downarrow & \operatorname{Tor} _ {1} (\bar {p} _ {r}, 1) & \downarrow & \downarrow \\ \operatorname{Tor} _ {1} ^ {\mathbf {A} _ {0}} (\mathbf {A} _ {0} / t ^ {r}, \mathbf {M}) & \longrightarrow & \operatorname{Tor} _ {1} ^ {\mathbf {A} _ {0}} (\overline {{\mathbf {A}}} _ {0} / \overline {{t}} ^ {r}, \mathbf {M}) & \longrightarrow & \mathbf {M} / x ^ {r - 1}   \mathbf {M} \xrightarrow {x _ {1}} \mathbf {M} / x ^ {r}   \mathbf {M}. \end{array}
\]

Now multiplication by \(x_{1}\) defines an injection of \(\mathbf{M}/\mathfrak{x}^{r}\mathbf{M}\) to \(\mathbf{M}/\mathfrak{x}^{r+1}\mathbf{M}\) : this follows immediately, by induction on \(r\) , from the exact sequence

\[
0 \rightarrow x ^ {r - 1} M / x ^ {r} M \rightarrow M / x ^ {r} M \rightarrow M / x ^ {r - 1} M \rightarrow 0
\]

and from the injectivity of the homothety with ratio \(x_{1}\) to \(\bigoplus_{r\geqslant 0}(\mathfrak{x}^{r}\mathrm{M}/\mathfrak{x}^{r+1}\mathrm{M})\) . Condition (v) therefore implies that the homomorphism \(\operatorname{Tor}_{1}^{\mathbf{A}_{0}}(\overline{p}_{r},1_{\mathbf{M}})\) is surjective for all \(r\geqslant 1\) . By composition with the isomorphisms \((\varphi_{\overline{\mathbf{A}}_{0}/\overline{\mathbf{t}}^{r+1}})^{-1}\) and \(\varphi_{\overline{\mathbf{A}}_{0}/\overline{\mathbf{t}}^{r}}\) , we deduce that the homomorphism

\[
\operatorname{Tor} _ {1} ^ {\overline {{\mathbf {A}}} _ {0}} \left(\bar {p} _ {r}, 1 _ {\mathbf {M}}\right): \operatorname{Tor} _ {1} ^ {\overline {{\mathbf {A}}} _ {0}} \left(\overline {{\mathbf {A}}} _ {0} / \bar {t} ^ {r + 1}, \mathbf {M}\right)\rightarrow \operatorname{Tor} _ {1} ^ {\overline {{\mathbf {A}}} _ {0}} \left(\overline {{\mathbf {A}}} _ {0} / \bar {t} ^ {r}, \mathbf {M}\right)
\]

is surjective for \(r \geqslant 1\) , whence condition (v') for \(\overline{\mathbf{M}}\) . This completes the proof of the theorem.

\subsection{Extension class associated with a regular sequence}

Let A be a ring, M an A-module, \(x = (x_{1}, \ldots, x_{n})\) a sequence of elements of A. Denote by \(M_{i}\) the A-module \(\mathbf{M}/(x_{1}\mathbf{M} + \cdots + x_{i-1}\mathbf{M})\) for \(i = 0, \ldots, n + 1\) , so that \(M_{0}\) and \(M_{1}\) identify with M and that \(\mathbf{M}_{n+1} = \mathbf{M}/(\mathbf{x})\mathbf{M}\) . Let

\[
\overline {{x}} _ {i}: \mathbf {M} _ {i - 1} \to \mathbf {M} _ {i}, \qquad i = 1, \dots , n,
\]

denote the A-homomorphism composed of the homothety of \(M_{i-1}\) of ratio \(x_{i}\) and of the canonical projection of \(M_{i-1}\) over \(M_{i}\) . Finally, let us denote by \(p: M_{n} \to M/(x)M\) the canonical projection. The diagram

\[
0 \longrightarrow \mathrm{M} \xrightarrow {\overline {{{x}}} _ {1}} \mathrm{M} _ {1} \xrightarrow {\overline {{{x}}} _ {2}} \mathrm{M} _ {2} \longrightarrow \dots \xrightarrow {\overline {{{x}}} _ {n}} \mathrm{M} _ {n} \xrightarrow {p} \mathrm{M} / (\boldsymbol {x}) \mathrm{M} \longrightarrow 0\tag{19}
\]

is an exact sequence if and only if the sequence x is M-regular. Suppose henceforth that the sequence x is regular for M. The element \(\theta_{x} \in \operatorname{Ext}_{A}^{n}(\mathbf{M}/(\mathbf{x})\mathbf{M}, \mathbf{M})\) associated with the exact sequence (19) is also said to be associated with the M-regular sequence x.

Let i be an integer, \(1 \leqslant i \leqslant n\) . Note that the sequence (19) can be decomposed into the two exact sequences

\begin{enumerate}
\def\labelenumi{(\arabic{enumi})}
\setcounter{enumi}{19}
\tightlist
\item
\end{enumerate}

\[
0 \longrightarrow \mathrm{M} \xrightarrow {\overline {{x}} _ {1}} \mathrm{M} _ {1} \xrightarrow {\overline {{x}} _ {2}} \mathrm{M} _ {2} \longrightarrow \dots \xrightarrow {\overline {{x}} _ {i}} \mathrm{M} _ {i} \longrightarrow \mathrm{M} / (x _ {1} \mathrm{M} + \dots + x _ {i} \mathrm{M}) \longrightarrow 0\tag{21}
\]

\[
0 \longrightarrow \mathrm{M} / (x _ {1} \mathrm{M} + \dots + x _ {i} \mathrm{M}) \xrightarrow {\overline {{x}} _ {i + 1}} \mathrm{M} _ {i + 1} \longrightarrow \dots
\]

\[
\longrightarrow \mathrm{M} _ {n} \stackrel {p} {\longrightarrow} \mathrm{M} / (x) \mathrm{M} \longrightarrow 0
\]

which are none other than the exact sequences associated with the sequence \((x_{1}, \ldots, x_{i})\) which is regular for M and to the sequence \((x_{i+1}, \ldots, x_{n})\) which is regular for \(\mathrm{M}/(x_{1} \mathrm{~M} + \cdots + x_{i} \mathrm{~M})\) . Denoting

\[
\theta_ {(x _ {1}, \dots , x _ {i})} \in \operatorname{Ext} _ {\mathbf {A}} ^ {i} (\mathbf {M} / (x _ {1} \mathbf {M} + \dots + x _ {i} \mathbf {M}), \mathbf {M})
\]

\[
\theta_ {(x _ {i + 1}, \dots , x _ {n})} \in \operatorname{Ext} _ {\mathbf {A}} ^ {n - i} (\mathbf {M} / (\boldsymbol {x}) \mathbf {M}, \mathbf {M} / (x _ {1} \mathbf {M} + \dots + x _ {i} \mathbf {M}))
\]

the extension classes associated with (20) and (21), we have, according to X, p.~118, prop. 3

\[
\theta_ {(x _ {1}, \dots , x _ {n})} = \theta_ {(x _ {1}, \dots , x _ {i})} \circ \theta_ {(x _ {i + 1}, \dots , x _ {n})}.\tag{22}
\]

Moreover, by prop. 5 (X, p.~157), the Koszul complex K.(x, M) is acyclic in degrees \(\neq n\) , whence an exact sequence

\[
0 \longrightarrow \mathrm{M} \xrightarrow {\partial^ {0}} \mathbf {K} ^ {1} (\boldsymbol {x}, \mathrm{M}) \xrightarrow {\partial^ {1}} \mathbf {K} ^ {2} (\boldsymbol {x}, \mathrm{M}) \xrightarrow {\partial^ {2}} \dots \longrightarrow \mathbf {K} ^ {n} (\boldsymbol {x}, \mathrm{M}) \xrightarrow {q} \mathrm{M} / (\boldsymbol {x}) \mathrm{M} \longrightarrow 0\tag{23}
\]

where we have identified \(\mathbf{K}^{0}(\boldsymbol{x},\mathbf{M})\) to \(\mathbf{M}\) and where \(q\) maps each element \(\boldsymbol{m}\in\mathbf{K}^{n}(\boldsymbol{x},\mathbf{M})\) on the class in \(\mathbf{M}/(\boldsymbol{x})\mathbf{M}\) of \(\boldsymbol{m}(1,2,\ldots,n)\in\mathbf{M}\) .

PROPOSITION 8. --- Suppose the sequence x is regular for M. The element of Ext \(_{A}^{n}\) (M/(x) M, M) associated with the exact sequence (23) is (-1) \(^{n(n+1)/2}\) . θ \(_{x}\) .

For \(i = 0,1,\dots,n\) let us define an A-linear map

\[
a ^ {i}: \mathbf {K} ^ {i} (\boldsymbol {x}, \mathrm{M}) \rightarrow \mathrm{M} _ {i} = \mathrm{M} / (x _ {1} \mathrm{M} + \dots + x _ {i - 1} \mathrm{M})
\]

as follows: if \(\boldsymbol{m} \in \mathbf{K}^i(\boldsymbol{x}, \mathbf{M})\) , \(a^i(\boldsymbol{m})\) is the class in \(\mathbf{M}_i\) of the element \(\boldsymbol{m}(1, 2, \ldots, i)\) of M. It is clear that \(a^0\) is the identity map of M and that \(p \circ a^n = q\) . Moreover \(a^{i+1} \circ \partial^i(\boldsymbol{m})\) is the image in \(\mathbf{M}_{i+1}\) of the element

\[
\sum_ {k = 1} ^ {i + 1} (- 1) ^ {k + 1} x _ {k} \boldsymbol {m} (1, 2, \dots , k - 1, k + 1, \dots , i + 1).
\]

Since \(x_{k}\) annihilates \(\mathbf{M}_{i+1}\) for \(k=1,\ldots,i,a^{i+1}\circ\partial^{i}(\boldsymbol{m})\) is also the image of

\[
(- 1) ^ {i} x _ {i + 1} \boldsymbol {m} (1, 2, \dots , i),
\]

hence

\[
a ^ {i + 1} \circ \partial^ {i} = (- 1) ^ {i} \overline {{{x}}} _ {i + 1} \circ a ^ {i}.
\]

By X, p.~120, cor. 1 and 2, the element of Ext \(_{A}^{n}\) (M/(x) M, M) associated with (23) is equal to \(\prod_{i=1}^{n} (-1)^{i} \cdot \theta_{x}\) , whence the assertion.

COROLLARY. — Suppose moreover that the modules \(\mathbf{M}/(x_{1}\mathbf{M}+\cdots+x_{i-1}\mathbf{M})\) are separated for the (x)-adic topology, and let \((a_{ij})\in\mathbf{GL}_{n}(\mathbf{A})\) . Set

\[
y _ {i} = \sum_ {j} a _ {i j} x _ {j} \text { and } \mathbf {y} = (y _ {1}, \dots , y _ {n}).
\]

Then the sequence y is regular for M, and one has \(\theta_{y} = \det (a_{ij})^{-1}\theta_{x}\) .

Indeed, the sequence y is regular for M by prop. 6 (X, p.~158) and theorem 1 (X, p.~160); the last assertion follows from prop. 8, and from prop. 4 of X, p.~119.

PROPOSITION 9. — Suppose the sequence x is regular for M. If N is an A-module such that (x) N = 0, one has Ext \(_{A}^{i}\) (N, M) = 0 for i \textless{} n, and the map α → θ \(_{x}\) ∘ α from Hom \(_{A}\) (N, M/(x) M) to Ext \(_{A}^{n}\) (N, M) (which is also the iterated connecting homomorphism associated with (19), cf.~X, p.~127, cor. 3) is bijective.

It remains to prove that the homomorphism \(\psi^{i}:\alpha\mapsto\theta_{x}\circ\alpha\) of \(\operatorname{Ext}_{\mathbf{A}}^{i-n}(\mathbf{N},\mathbf{M}/(\mathbf{x})\mathbf{M})\) to \(\operatorname{Ext}_{\mathbf{A}}^{i}(\mathbf{N},\mathbf{M})\) is bijective for \(i\leqslant n\) . Let us reason by induction on \(n\) , the assertion being trivial for \(n=0\) . Set \(\mathbf{M}_{1}=\mathbf{M}/x_{1}\mathbf{M}\) , \(\mathbf{x}^{\prime}=(x_{2},...,x_{n})\) , so that \(\mathbf{x}^{\prime}\) is \(\mathbf{M}_{1}\) -regular. By the induction hypothesis, the homomorphism

\[
\overline {{{{\psi}}}} ^ {i - 1}: \alpha \mapsto \theta_ {\boldsymbol {x} ^ {\prime}} \circ \alpha
\]

of \(\operatorname{Ext}_{\mathbf{A}}^{i-n}(\mathbf{N}, \mathbf{M}/(\mathbf{x})\mathbf{M})\) to \(\operatorname{Ext}_{\mathbf{A}}^{i-1}(\mathbf{N}, \mathbf{M}_1)\) is bijective for \(i < n\) . Moreover, consider the exact sequence

\[
0 \longrightarrow \mathrm{M} \xrightarrow {(x _ {1}) _ {\mathrm{M}}} \mathrm{M} \longrightarrow \mathrm{M} _ {1} \longrightarrow 0;
\]

the connecting homomorphism \(\operatorname{Ext}_{\mathbf{A}}^{i}(\mathbf{N}, \mathbf{M}_{1}) \to \operatorname{Ext}_{\mathbf{A}}^{i+1}(\mathbf{N}, \mathbf{M})\) associated is \(\varphi^{i}: \beta \mapsto \theta_{x_{1}} \circ \beta\) (X, p.~125, prop. 5);

since we have \(\operatorname{Ext}^i(1_{\mathbf{N}}, (x_1)_{\mathbf{M}}) = \operatorname{Ext}^i((x_1)_{\mathbf{N}}, 1_{\mathbf{M}}) = 0\) , we deduce exact sequences

\[
0 \longrightarrow \operatorname{Ext} _ {\mathrm{A}} ^ {i} (\mathrm{N}, \mathrm{M}) \longrightarrow \operatorname{Ext} _ {\mathrm{A}} ^ {i} (\mathrm{N}, \mathrm{M} _ {1}) \xrightarrow {\varphi^ {i}} \operatorname{Ext} _ {\mathrm{A}} ^ {i + 1} (\mathrm{N}, \mathrm{M}) \longrightarrow 0.
\]

Since \(\operatorname{Ext}_{\mathbf{A}}^{i}(\mathbf{N}, \mathbf{M}_{1}) = 0\) for \(i < n - 1\) , we deduce that \(\operatorname{Ext}_{\mathbf{A}}^{i + 1}(\mathbf{N}, \mathbf{M}) = 0\) for \(i < n - 1\) , that is, \(i + 1 < n\) ; it follows that \(\varphi^{i}\) is bijective for \(i < n\) . Since \(\psi^{i}(\alpha) = \theta_{x} \circ \alpha = \theta_{x_{1}} \circ \theta_{x'} \circ \alpha = \varphi^{i - 1} \circ \overline{\psi}^{i}(\alpha)\) for \(\alpha \in \operatorname{Ext}_{\mathbf{A}}^{i - n}(\mathbf{N}, \mathbf{M}/(\mathbf{x})\mathbf{M})\) , \(\psi^{i}\) is indeed bijective for \(i \leqslant n\) .

\backmatter
\specialchapter{Exercises}
